\ifdefined\HAFOGeoPosteriorOnly\else
\section{Propósito y alcance geométrico--topológico}
\label{sec:geo-proposito}

Los capítulos anteriores responden principalmente preguntas algebraicas y
numéricas: cómo realizar un sistema en forma de Lur'e, cómo producir una
semilla, cómo continuarla y qué diagnósticos calcular. La teoría geométrica y
topológica cambia el centro de atención. En lugar de preguntar solamente
«¿cuál es la trayectoria?», pregunta:

\begin{itemize}
  \item ¿sobre qué espacio vive el sistema y qué estructura de ese espacio
        debe respetarse?;
  \item ¿qué conjuntos transporta el flujo en sí mismos?;
  \item ¿qué propiedades sobreviven a un cambio de coordenadas o a una
        reparametrización del tiempo?;
  \item ¿cómo están organizadas las cuencas, sus fronteras y las separatrices?;
  \item ¿qué mecanismo geométrico podría sostener la periodicidad o el caos?;
  \item ¿qué parte de una conclusión procede de un teorema y qué parte es sólo
        evidencia computacional finita?
\end{itemize}

\ifHAFOPedagogy
\subsection{Objetivos de aprendizaje}

Al terminar esta guía, el lector debería poder:

\begin{enumerate}
  \item identificar correctamente el espacio de fases de Chua, MAVPD,
        Kalman--Fitts y PLL;
  \item definir flujo, órbita, invariancia, conjunto \(\omega\)-límite,
        atractor, cuenca y frontera de cuenca;
  \item distinguir conjugación topológica de equivalencia orbital;
  \item usar la linealización y los teoremas de Hartman--Grobman y de la
        variedad estable sólo cuando se cumplen sus hipótesis;
  \item construir una sección de Poincaré y relacionar su derivada con los
        multiplicadores de Floquet;
  \item reconocer qué faltaría para convertir una nube de retornos en una
        prueba de homoclinicidad o de herradura;
  \item usar índice y grado topológico como herramientas de existencia, sin
        confundir existencia con atracción u ocultedad;
  \item redactar cada conclusión separando visualización, evidencia numérica y
        demostración.
\end{enumerate}

\subsection{Lectura previa exacta}

Antes de continuar, complete las etapas 1--5 de
\cref{sec:programa-estudio-bibliografico}. Allí se indican páginas impresas de
los cuatro libros físicos, páginas del visor de los PDF locales, ejercicios y
una comprobación para decidir si puede avanzar. El orden mínimo es
\[
\text{campo y flujo}
\longrightarrow \text{espacio de fases}
\longrightarrow \text{Poincaré--Floquet}
\longrightarrow \text{grado}
\longrightarrow \text{caos y cuencas}.
\]
No empiece por conexiones, curvatura o geometría fractal: esas lecturas son de
profundización y aparecen al final de la ruta.
\fi

\subsection{Escalera de afirmaciones}
\label{geo:subsec:escalera}

Para evitar que una gráfica cargue más significado del que posee, se usará la
siguiente escalera.

\begin{table}[H]
\centering
\caption{Niveles de una afirmación geométrica o topológica.}
\label{geo:tab:niveles-evidencia}
\small
\begin{tabularx}{\textwidth}{@{}L{0.12\textwidth}L{0.22\textwidth}YY@{}}
\toprule
Nivel & Producto & Qué permite decir & Qué no permite decir \\
\midrule
G0 & Visualización & Sugiere una órbita, una separación o una estructura que
merece formularse. & No prueba invariancia, convergencia, caos ni ocultedad. \\
G1 & Evidencia numérica reproducible & Apoya una propiedad dentro de un
horizonte, tolerancias, dominio y conjunto de semillas declarados. & No cubre
tiempo infinito ni un continuo de condiciones iniciales. \\
G2 & Teorema local con hipótesis verificadas & Demuestra una conclusión local:
hiperbolicidad, dimensión de variedades o estabilidad de un ciclo, por
ejemplo. & No clasifica automáticamente la cuenca global. \\
G3 & Argumento global o prueba asistida & Certifica una propiedad en un dominio
mediante grado, regiones atrapantes, cotas rigurosas o intersecciones
transversales encerradas. & No debe extenderse fuera del dominio certificado.
\\
\bottomrule
\end{tabularx}
\end{table}

Un sistema de álgebra computacional puede establecer exactamente una identidad
algebraica y ayudar a verificar una hipótesis de G2 o G3. Una integración
numérica sigue siendo G1 mientras no incluya cotas rigurosas que cubran el
continuo relevante.

\section{Lenguaje mínimo de topología y geometría diferencial}
\label{sec:geo-lenguaje}

Esta parte se restringe primero a ecuaciones autónomas de orden entero
\begin{equation}
  \dot x=f(x),\qquad x\in M,
  \label{geo:eq:edo-manifold}
\end{equation}
donde \(M\) es una variedad diferenciable, usualmente \(\R^n\) o el cilindro
\(\R\times S^1\). Esta restricción no es decorativa: el flujo y la propiedad de
grupo que se construirán aquí no se trasladan sin más a un PVI causal de
Caputo con memoria. La extensión fraccionaria debe formularse por separado.
El lenguaje topológico y diferencial de esta parte sigue
\cite{NahmadAchar2022Topologia}; la teoría de flujo y estabilidad se apoya en
\cite{JaurezEtAl2021TGEDI,JaurezEtAl2021TGEDII,HirschSmaleDevaney2013}.

\subsection{Topología: proximidad sin fijar una regla de medida}

\begin{definition}[Espacio topológico]
Un espacio topológico es un par \((M,\mathcal T)\), donde \(\mathcal T\) es una
familia de subconjuntos llamados \emph{abiertos} que contiene a
\(\varnothing\) y \(M\), es cerrada bajo uniones arbitrarias y bajo
intersecciones finitas. Una \emph{vecindad} de \(x\) es un conjunto que
contiene un abierto que contiene a \(x\).
\end{definition}

\begin{definition}[Base, separación, compacidad y conexidad]
Una familia \(\mathcal B\subset\mathcal T\) es una \emph{base} si todo abierto
es unión de elementos de \(\mathcal B\); una subbase genera una base mediante
intersecciones finitas. El espacio es \emph{Hausdorff} si dos puntos distintos
poseen vecindades disjuntas. Es \emph{compacto} si toda cubierta abierta admite
una subcubierta finita, y es \emph{conexo} si sus únicos subconjuntos que son a
la vez abiertos y cerrados son \(\varnothing\) y el espacio completo.
\end{definition}

En espacios métricos la compacidad equivale a que toda sucesión posea una
subsucesión convergente; esa equivalencia no debe trasladarse sin hipótesis a
un espacio topológico arbitrario. Compacidad controla argumentos de
extracción, mientras conexidad impide separar el objeto en dos abiertos
disjuntos no vacíos.

\begin{definition}[Producto y cociente]
La topología producto en \(M\times N\) tiene como base los conjuntos
\(U\times V\), con \(U\subset M\) y \(V\subset N\) abiertos. Dada una relación
de equivalencia \(\sim\), el cociente \(M/{\sim}\) lleva la topología para la
cual un conjunto \(W\) es abierto si \(\pi^{-1}(W)\) es abierto, donde
\(\pi:M\to M/{\sim}\) es la proyección canónica.
\end{definition}

Así, \(S^1\simeq\R/(2\pi\mathbb Z)\) identifica ángulos que difieren por una
vuelta completa, y el cilindro del PLL usa la topología producto
\(\R\times S^1\).
\ifHAFOPedagogy
Como ejercicio de dominio, pruebe que la proyección
\(\theta\mapsto(\cos\theta,\sin\theta)\) induce un homeomorfismo del cociente
sobre el círculo; señale dónde usa compacidad y la propiedad de Hausdorff.
\fi

Una métrica \(d\) produce una topología mediante sus bolas abiertas, pero la
topología conserva menos información que la métrica. Continuidad, compacidad,
conexidad y frontera son nociones topológicas; longitudes, ángulos y radios
numéricos requieren estructura adicional.

\begin{definition}[Continuidad y homeomorfismo]
Una función \(h:M\to N\) es continua si la preimagen de todo abierto de \(N\)
es abierta en \(M\). Es un \emph{homeomorfismo} si es biyectiva y tanto \(h\)
como \(h^{-1}\) son continuas.
\end{definition}

Un homeomorfismo puede estirar una bola hasta convertirla en un elipsoide, sin
romper conexidad ni vecindades. Ésta es la razón por la cual la ocultedad,
formulada mediante vecindades, puede transportarse topológicamente, mientras
que un contrato de «radio \(10^{-3}\)» no se transporta conservando el mismo
número.

\begin{definition}[Interior, clausura y frontera]
Para \(A\subset M\), su interior \(\operatorname{int}A\) es el mayor abierto
contenido en \(A\); su clausura \(\overline A\) es el menor cerrado que lo
contiene; y
\begin{equation}
  \partial A=\overline A\setminus\operatorname{int}A
  \label{geo:eq:boundary}
\end{equation}
es su frontera.
\end{definition}

La frontera de una cuenca no es la última curva coloreada por una malla. Es el
conjunto definido por~\eqref{geo:eq:boundary}; cualquier vecindad de uno de sus
puntos encuentra tanto puntos de la cuenca como puntos fuera de ella.

\subsection{Variedades, vectores tangentes y campos}

\begin{definition}[Variedad diferenciable y campo vectorial]
Una variedad diferenciable \(M\) de dimensión \(n\) es un espacio que
localmente posee coordenadas en abiertos de \(\R^n\), con cambios de carta
diferenciables. El espacio tangente en \(x\) se denota \(T_xM\). Un campo
vectorial es una asignación \(f(x)\in T_xM\) que varía con la regularidad
declarada.
\end{definition}

En \(\R^n\) se identifica \(T_x\R^n\simeq\R^n\), por lo que el vector
\(f(x)\) se dibuja con origen en \(x\). En el PLL, en cambio,
\begin{equation}
  M_{\mathrm{PLL}}=\R\times S^1,
  \label{geo:eq:pll-phase-space}
\end{equation}
y los pares \((x,\theta)\) y \((x,\theta+2\pi k)\) representan el mismo estado.
El rectángulo dibujado en el plano es sólo una carta cortada y pegada por sus
bordes angulares.

Una distancia compatible útil para visualización y comparación numérica es
\begin{equation}
 d_{\mathrm{cil}}((x,\theta),(\tilde x,\tilde\theta))
 =\sqrt{(x-\tilde x)^2+
 \operatorname{wrap}_{-\pi\le\,\cdot\,<\pi}(\theta-\tilde\theta)^2}.
 \label{geo:eq:cylindrical-distance}
\end{equation}
No es legítimo usar \(|\theta-\tilde\theta|\) sin envolver y después atribuir
significado topológico al salto artificial en \(\pm\pi\).

\paragraph{De cartas a formas.}
Una carta es un homeomorfismo \(\phi:U\to\phi(U)\subset\R^n\); un atlas de
clase \(C^r\) exige que los cambios \(\psi\circ\phi^{-1}\) sean \(C^r\). Una
subvariedad se reconoce localmente por coordenadas en las que queda descrita
por la anulación de algunas componentes. El cotangente \(T_x^*M\) es el dual
de \(T_xM\), y una \(k\)-forma es una sección de \(\Lambda^kT^*M\). El
producto exterior, la derivada exterior y el pull-back permiten formular
integración y Stokes sin depender de una carta; una orientación selecciona
atlas cuyas transiciones conservan orientación
\cite{NahmadAchar2022Topologia}. La subsección sobre formas retomará estos
objetos para convertir divergencia en balance de flujo.

\ifHAFOPedagogy
\paragraph{Ejercicio de enlace.}
Para \(h(r,\theta)=(r\cos\theta,r\sin\theta)\), calcule
\(h^*(\dd x\wedge\dd y)\), identifique dónde deja de ser una carta y explique
qué información de orientación contiene el signo del jacobiano.
\fi

\subsection{Cambio de coordenadas como transporte de un campo}

Sea \(h:M\to N\) un difeomorfismo y \(y=h(x)\). Por la regla de la cadena,
el campo transformado es
\begin{equation}
  g(y)=D h(h^{-1}(y))f(h^{-1}(y)).
  \label{geo:eq:pushforward}
\end{equation}
Esta ecuación es el \emph{empuje hacia delante} del campo, \(g=h_*f\). No
basta cambiar los nombres de las variables: el diferencial \(Dh\) es parte de
la transformación.

\begin{proposition}[Transporte de curvas integrales]
Supóngase que \(h\) es un difeomorfismo y que \(x(t)\) resuelve
\(\dot x=f(x)\). Entonces \(y(t)=h(x(t))\) resuelve
\(\dot y=g(y)\), con \(g\) dado por~\eqref{geo:eq:pushforward}.
\end{proposition}

\begin{proof}
La regla de la cadena da
\[
 \dot y(t)=Dh(x(t))\dot x(t)=Dh(x(t))f(x(t))=g(h(x(t))).
\]
Como \(y(t)=h(x(t))\), se obtiene la ecuación afirmada.
\end{proof}

Para una transformación afín \(x=Tz+c\), una bola física alrededor de
\(x_e\) se convierte en
\begin{equation}
 (z-z_e)^{\mathsf T}T^{\mathsf T}T(z-z_e)\le r^2.
 \label{geo:eq:ball-to-ellipsoid}
\end{equation}
Así se conecta la equivalencia afín ya usada por las realizaciones D0/T1 con
la geometría de los contratos de cuenca.

\subsection{Regularidad de los cuatro casos enteros}

\begin{table}[H]
\centering
\caption{Espacio de fases y regularidad relevante.}
\label{geo:tab:regularity-cases}
\small
\begin{tabularx}{\textwidth}{@{}L{0.18\textwidth}L{0.16\textwidth}L{0.25\textwidth}Y@{}}
\toprule
Caso & Espacio & Regularidad & Precaución geométrica \\
\midrule
Chua PWL & \(\R^3\) & Campo continuo y globalmente Lipschitz; suave por
regiones & \(Df\) salta en \(x=\pm1\); los teoremas \(C^1\) deben aplicarse por
regiones o con una versión no suave. \\
MAVPD & \(\R^3\) & Campo polinómico \(C^\infty\) & La existencia global no se
deduce de que el campo sea polinómico; se requieren cotas de no explosión o
una región atrapante para la órbita considerada. \\
Kalman--Fitts & \(\R^4\) & Campo \(C^\infty\), con no linealidad
\(\tanh\) globalmente Lipschitz & Una sección de Poincaré es tridimensional;
un retrato en dos coordenadas es sólo una proyección. \\
PLL lead--lag & \(\R\times S^1\) & Campo \(C^\infty\) en el cilindro & Deben
envolverse ángulos, retornos y vecindades; una órbita corredora es cerrada en
el cilindro aunque no lo parezca en el recubrimiento \(\R^2\). \\
\bottomrule
\end{tabularx}
\end{table}

\section{Del problema de valor inicial al flujo}
\label{sec:geo-flujo}

\subsection{Existencia, unicidad y dependencia de la condición inicial}

\begin{proposition}[Picard--Lindelöf local, forma usada aquí]
Sea \(f:U\subset\R^n\to\R^n\) localmente Lipschitz. Para cada
\(x_0\in U\) existe una única solución maximal
\(x(\,\cdot\,;x_0)\) de~\eqref{geo:eq:edo-manifold} definida en un intervalo
abierto que contiene a \(0\). La solución depende continuamente de \(x_0\) en
cualquier rectángulo compacto donde permanezcan las soluciones.
\end{proposition}

La hipótesis de unicidad es el fundamento de la geometría de órbitas: dos
trayectorias distintas no pueden cruzarse en el mismo instante y estado. La
existencia local no implica existencia para todo \(t\in\R\). Un campo con
crecimiento a lo sumo lineal sí proporciona una vía estándar hacia completitud
global; para campos polinómicos de grado mayor se necesitan cotas adicionales.

\begin{definition}[Flujo local y flujo completo]
El flujo local asociado con~\eqref{geo:eq:edo-manifold} es
\begin{equation}
 \varphi_t(x_0)=x(t;x_0),
 \label{geo:eq:flow}
\end{equation}
en el dominio donde existe la solución. Es un flujo completo si está definido
para todo \((t,x)\in\R\times M\) y satisface
\begin{equation}
 \varphi_0=\operatorname{id},\qquad
 \varphi_{t+s}=\varphi_t\circ\varphi_s,
 \qquad \varphi_{-t}=\varphi_t^{-1}.
 \label{geo:eq:flow-group}
\end{equation}
\end{definition}

\begin{proposition}[La unicidad produce la propiedad de flujo]
Si las dos composiciones están definidas, entonces
\(\varphi_{t+s}(x)=\varphi_t(\varphi_s(x))\).
\end{proposition}

\begin{proof}
Las curvas \(u(t)=\varphi_{t+s}(x)\) y
\(v(t)=\varphi_t(\varphi_s(x))\) resuelven la misma ecuación autónoma y
satisfacen \(u(0)=v(0)=\varphi_s(x)\). La unicidad implica \(u=v\).
\end{proof}

\begin{remark}[Por qué autonomía importa]
En una ecuación no autónoma \(\dot x=f(t,x)\) se usa un proceso de dos tiempos
\(\varphi(t,t_0,x_0)\). La identidad~\eqref{geo:eq:flow-group} no puede
escribirse eliminando \(t_0\) a menos que se amplíe el espacio con una variable
de fase temporal.
\end{remark}

\subsection{Órbitas y conjuntos invariantes}

\begin{definition}[Órbita, equilibrio y órbita periódica]
La órbita positiva, negativa y completa de \(x\) son
\begin{align}
 \mathcal O^+(x)&=\{\varphi_t(x):t\ge0\},\nonumber\\
 \mathcal O^-(x)&=\{\varphi_t(x):t\le0\},\nonumber\\
 \mathcal O(x)&=\{\varphi_t(x):t\in\R\}.
 \label{geo:eq:orbits}
\end{align}
Un punto \(e\) es un equilibrio si \(f(e)=0\), equivalentemente
\(\varphi_t(e)=e\). Una órbita es periódica si existe \(T>0\) tal que
\(\varphi_T(x)=x\); el menor \(T>0\) es su periodo mínimo.
\end{definition}

\begin{definition}[Invariancia]
Un conjunto \(A\subset M\) es positivamente invariante si
\(\varphi_t(A)\subset A\) para \(t\ge0\); negativamente invariante si ocurre
para \(t\le0\); e invariante si \(\varphi_t(A)=A\) para todo \(t\in\R\).
\end{definition}

Un retrato de fase dentro de una caja durante tiempo finito no demuestra que
la caja sea invariante. Para un dominio suave sí existe una prueba elemental.

\begin{proposition}[Criterio de entrada estricta]
Sea \(N=\{x:h(x)\le0\}\) compacto, con \(h\in C^1\) y
\(\nabla h\ne0\) sobre \(\partial N\). Si
\begin{equation}
 \nabla h(x)\cdot f(x)<0
 \quad\text{para todo }x\in\partial N,
 \label{geo:eq:strict-inward}
\end{equation}
y las soluciones son únicas, entonces \(N\) es positivamente invariante.
\end{proposition}

\begin{proof}
Si una solución que empieza en \(N\) saliera por primera vez en \(t_*\),
entonces \(h(x(t_*))=0\) y, para cruzar hacia \(h>0\), su derivada derecha
tendría que ser no negativa. Pero
\(\frac{d}{dt}h(x(t_*))=\nabla h(x(t_*))\cdot f(x(t_*))<0\), contradicción.
\end{proof}

\begin{center}
\resizebox{0.96\linewidth}{!}{%
\begin{tikzpicture}[
  >=Latex,
  every node/.style={font=\small},
  flow/.style={-{Latex[length=2.4mm]},thick,draw=Azul},
  vector/.style={-{Latex[length=2.4mm]},very thick},
  point/.style={circle,fill=GrisOscuro,inner sep=1.8pt},
  panel/.style={draw=GrisOscuro,rounded corners=5pt,thick}
]
  % Propiedad de composición del flujo.
  \draw[panel] (0,0) rectangle (6.25,5.05);
  \node[font=\bfseries] at (3.12,4.68) {Composición temporal};
  \node[point,label=left:$x$] (x) at (0.95,3.25) {};
  \node[point,label=above:$\varphi_s(x)$] (xs) at (4.95,3.25) {};
  \node[point,label=below:$\varphi_{t+s}(x)$] (xts) at (4.95,0.95) {};
  \draw[flow] (x) -- node[above] {$\varphi_s$} (xs);
  \draw[flow] (xs) -- node[right] {$\varphi_t$} (xts);
  \draw[flow,dashed] (x) -- node[below left] {$\varphi_{t+s}$} (xts);
  \node[align=center,text=GrisOscuro] at (3.0,0.38)
    {$\varphi_t\!\circ\!\varphi_s=\varphi_{t+s}$\\por unicidad};

  % Criterio geométrico de entrada estricta.
  \draw[panel] (6.65,0) rectangle (13.9,5.05);
  \node[font=\bfseries] at (10.27,4.68) {Invariancia positiva};
  \path[fill=AzulClaro,draw=Azul,very thick]
    plot[smooth cycle,tension=0.85] coordinates {
      (7.45,1.25) (8.05,3.65) (10.05,4.05)
      (12.35,3.60) (13.0,1.55) (10.35,0.65)
    };
  \node[font=\large] at (9.75,2.25) {$N=\{h\leq0\}$};
  \node[point,label=above left:$x_b\in\partial N$] (xb) at (12.32,3.60) {};
  \draw[vector,draw=GrisOscuro] (xb) -- ++(1.08,0.65)
    node[right,align=left] {$n=\nabla h/\|\nabla h\|$\\normal exterior};
  \draw[vector,draw=Rojo] (xb) -- ++(-1.30,-0.68)
    node[below left] {$f(x_b)$};
  \draw[flow] (8.25,1.55) .. controls (9.15,1.05) and (10.55,1.10) .. (11.35,1.72);
  \node[fill=white,inner sep=2pt] at (10.0,0.35)
    {$\nabla h(x_b)\!\cdot\!f(x_b)<0\ \Longrightarrow\ \varphi_t(N)\subset N$};
\end{tikzpicture}%
}

\smallskip
\begin{minipage}{0.92\linewidth}
\small\textbf{Lectura pedagógica.}
La identidad de composición convierte las soluciones únicas en un flujo. El
signo estrictamente negativo de la derivada de $h$ en la frontera impide la
salida hacia $h>0$; prueba invariancia positiva, no necesariamente la igualdad
\(\varphi_t(N)=N\) para todo tiempo.
\end{minipage}
\end{center}


La desigualdad no estricta requiere una formulación de tangencia o viabilidad
más cuidadosa; reemplazar \(<0\) por \(\le0\) sin revisar hipótesis puede
invalidar el argumento de primer cruce.

\subsection{Ejemplo resuelto: las superficies de conmutación de Chua}

En~\eqref{int:eq:chua}, las superficies de conmutación son
\(S_\pm=\{x=\pm1\}\). Con \(h_+(x,y,z)=x-1\) y
\(h_-(x,y,z)=x+1\), las derivadas a lo largo del campo son
\begin{align}
 \dot h_+\big|_{S_+}&=\alpha(y-1-m_0),\nonumber\\
 \dot h_-\big|_{S_-}&=\alpha(y+1+m_0).
 \label{geo:eq:chua-switch-lie}
\end{align}
Por tanto, el cruce es transversal salvo sobre las rectas de tangencia
\begin{equation}
 S_+\cap\{y=1+m_0\},\qquad
 S_-\cap\{y=-1-m_0\}.
 \label{geo:eq:chua-grazing-lines}
\end{equation}
Las tres regiones PWL no son, en general, invariantes: el signo de
\(\dot h_\pm\) depende de \(y\). Una figura que coloree cada región no debe
interpretarse como una partición en componentes dinámicamente cerradas.

\subsection{Simetría como propiedad del flujo}

\begin{proposition}[Simetría por una involución]
Sea \(R:M\to M\) un difeomorfismo con \(R^2=\operatorname{id}\). Si
\begin{equation}
 f(Rx)=DR(x)f(x),
 \label{geo:eq:equivariant-field}
\end{equation}
entonces \(R\varphi_t(x)=\varphi_t(Rx)\) siempre que ambas expresiones estén
definidas. En consecuencia, \(R\) transporta órbitas, conjuntos invariantes y
cuencas.
\end{proposition}

\begin{proof}
La curva \(R\varphi_t(x)\) tiene derivada
\(DR(\varphi_t(x))f(\varphi_t(x))=f(R\varphi_t(x))\) y empieza en \(Rx\).
Por unicidad coincide con \(\varphi_t(Rx)\).
\end{proof}

Para \(R(x)=-x\), la condición es simplemente \(f(-x)=-f(x)\). Chua PWL,
MAVPD y Kalman--Fitts poseen esta simetría impar. Por ejemplo, si
\(\mathcal A\) es un atractor MAVPD, entonces \(-\mathcal A\) también lo es y
\begin{equation}
 \mathcal B(-\mathcal A)=-\mathcal B(\mathcal A).
 \label{geo:eq:symmetric-basins}
\end{equation}
Los dos conjuntos pueden coincidir o formar un par distinto; la simetría por
sí sola no decide cuál alternativa ocurre. En el expediente MAVPD, los ciclos
internos aparecen como un par simétrico, mientras que la extensión externa
debe comprobarse por sus propias trayectorias.

\section{Conjuntos límite, atracción y cuencas}
\label{sec:geo-limit-basins}

\subsection{Qué queda después del transitorio}

\begin{definition}[Conjunto \texorpdfstring{\(\omega\)-límite}{omega-límite}]
Para una órbita positiva completa,
\begin{equation}
 \omega(x)=
 \bigcap_{s\ge0}\overline{\{\varphi_t(x):t\ge s\}}.
 \label{geo:eq:omega-limit}
\end{equation}
Equivalentemente, \(y\in\omega(x)\) si existe una sucesión
\(t_k\to+\infty\) tal que \(\varphi_{t_k}(x)\to y\). El conjunto
\(\alpha(x)\) se define análogamente con \(t_k\to-\infty\).
\end{definition}

\begin{proposition}[Estructura básica de \texorpdfstring{\(\omega(x)\)}{omega(x)}]
Sea \(\varphi\) un flujo continuo y completo. Si
\(\mathcal O^+(x)\) tiene clausura compacta, entonces \(\omega(x)\) es no
vacío, compacto, conexo e invariante.
\end{proposition}

\begin{proof}
Los conjuntos
\(K_s=\overline{\{\varphi_t(x):t\ge s\}}\) son compactos no vacíos, conexos y
anidados. Su intersección es no vacía, compacta y conexa. Si
\(y=\lim_k\varphi_{t_k}(x)\), la continuidad da
\(\varphi_\tau(y)=\lim_k\varphi_{t_k+\tau}(x)\in\omega(x)\). Usando
completitud y \(t_k-\tau\to\infty\) se obtiene también la inclusión inversa,
de modo que el conjunto es invariante.
\end{proof}

La precompacidad es una hipótesis, no una consecuencia de que la gráfica
quepa en una ventana. Una rutina que detiene la integración al llegar a un
umbral de divergencia tampoco construye un conjunto \(\omega\)-límite.

\subsection{Atrayente, atractor y cuenca}

Para un conjunto \(A\) en un espacio métrico, se escribe
\(\operatorname{dist}(x,A)=\inf_{a\in A}d(x,a)\).

\begin{definition}[Atracción y cuenca]
Un compacto invariante \(A\) atrae a \(x\) si
\begin{equation}
 \operatorname{dist}(\varphi_t(x),A)\longrightarrow0
 \quad(t\to\infty).
 \label{geo:eq:attraction}
\end{equation}
Su cuenca es
\begin{equation}
 \mathcal B(A)=\{x\in M:
 \operatorname{dist}(\varphi_t(x),A)\to0\}.
 \label{geo:eq:basin}
\end{equation}
En esta guía, un atractor asintóticamente estable es un compacto invariante
que atrae una vecindad y es estable en el sentido de Lyapunov.
\end{definition}

La estabilidad de Lyapunov significa: para cada vecindad \(U\) de \(A\)
existe una vecindad \(V\) de \(A\) tal que
\(\varphi_t(V)\subset U\) para todo \(t\ge0\). Atracción y estabilidad son
propiedades distintas; una definición de atractor debe declarar cuál de las
convenciones adopta.

\begin{proposition}[La cuenca de un atractor asintóticamente estable es abierta]
Bajo las hipótesis de la definición anterior y continuidad del flujo,
\(\mathcal B(A)\) es abierta.
\end{proposition}

\begin{proof}
Sea \(x\in\mathcal B(A)\). Elija una vecindad atractora \(V\) de \(A\) que,
por estabilidad, permanezca dentro de una vecindad donde la convergencia está
garantizada. Existe \(T\) con \(\varphi_T(x)\in V\). La continuidad de
\(\varphi_T\) produce una vecindad \(O\) de \(x\) tal que
\(\varphi_T(O)\subset V\). Todo punto de \(O\) converge a \(A\), por lo que
\(O\subset\mathcal B(A)\).
\end{proof}

\begin{definition}[Atractor oculto, formulación topológica]
Sea \(\mathcal E\) el conjunto de equilibrios. Un atractor \(A\) es oculto si,
para cada \(e\in\mathcal E\), existe una vecindad abierta \(U_e\) tal que
\begin{equation}
 U_e\cap\mathcal B(A)=\varnothing.
 \label{geo:eq:hidden-topological}
\end{equation}
Es autoexcitado si su cuenca intersecta toda vecindad suficientemente pequeña
de al menos un equilibrio en la convención adoptada.
\end{definition}

La condición~\eqref{geo:eq:hidden-topological} es cualitativa y global. Probar
cero contactos en 48, 108 o 504 lanzamientos es evidencia G1 de una versión
discretizada; no cuantifica el abierto completo \(U_e\).

\subsection{Fronteras de cuenca y separatrices}

\begin{proposition}[Invariancia de la frontera de cuenca]
Sea \(\varphi\) un flujo completo de homeomorfismos y sea \(A\) un atractor.
Entonces
\begin{equation}
 \varphi_t(\mathcal B(A))=\mathcal B(A),\qquad
 \varphi_t(\partial\mathcal B(A))=\partial\mathcal B(A)
 \quad(t\in\R).
 \label{geo:eq:basin-boundary-invariance}
\end{equation}
\end{proposition}

\begin{proof}
Desplazar una órbita una cantidad finita no cambia su límite cuando
\(t\to\infty\), de modo que la cuenca es invariante. Todo
\(\varphi_t\) es un homeomorfismo; los homeomorfismos transportan interior,
clausura y, por~\eqref{geo:eq:boundary}, frontera. Esto da la segunda igualdad.
\end{proof}

\begin{center}
\begin{minipage}{\linewidth}
\centering
\resizebox{0.96\linewidth}{!}{%
\begin{tikzpicture}[
  >=Latex,
  every node/.style={font=\small},
  flow/.style={-{Latex[length=2.3mm]},thick,draw=GrisOscuro},
  point/.style={circle,fill=GrisOscuro,inner sep=1.8pt}
]
  \begin{scope}
    \clip[rounded corners=7pt] (0.2,0.25) rectangle (13.8,6.0);
    \fill[VerdeClaro] (0.2,0.25) rectangle (13.8,6.0);
    \fill[AzulClaro]
      (0.2,0.25) -- (6.35,0.25)
      .. controls (5.55,1.15) and (6.85,1.75) .. (6.05,2.55)
      .. controls (5.30,3.25) and (6.65,4.25) .. (6.35,6.0)
      -- (0.2,6.0) -- cycle;
  \end{scope}
  \draw[rounded corners=7pt,very thick,GrisOscuro] (0.2,0.25) rectangle (13.8,6.0);
  \node[anchor=north west,font=\bfseries] at (0.45,5.82) {Espacio de fases \(M\)};
  \node[font=\bfseries,Azul] at (2.55,5.15) {$\mathcal B(A_1)$};
  \node[font=\bfseries,Verde] at (10.55,5.15) {$\mathcal B(A_h)$};

  % Frontera común, punto silla y transporte sobre la frontera.
  \draw[very thick,Naranja]
    (6.35,0.25)
    .. controls (5.55,1.15) and (6.85,1.75) .. (6.05,2.55)
    .. controls (5.30,3.25) and (6.65,4.25) .. (6.35,6.0);
  \draw[very thick,GrisOscuro,densely dashed]
    (6.00,2.30) .. controls (5.55,3.05) and (6.55,3.75) .. (6.24,4.45);
  \node[point,label=right:$S$] (s) at (6.00,3.05) {};
  \node[align=left,fill=white,inner sep=2pt] at (7.78,4.15)
    {componente candidata\\$W^s(S)$};
  \node[point,label=left:$x$] (xb0) at (6.20,1.15) {};
  \node[point,label=right:$\varphi_t(x)$] (xb1) at (6.08,2.35) {};
  \draw[flow] (xb0) .. controls (5.78,1.62) and (6.35,1.95) .. (xb1);
  \node[rotate=90,Naranja,font=\bfseries] at (5.45,4.90)
    {$\partial\mathcal B(A_h)$};

  % Atractores y vecindad de equilibrio que no toca la cuenca oculta.
  \draw[very thick,Azul] (2.55,2.55) circle (0.43);
  \fill[Azul] (2.55,2.55) circle (2pt);
  \node[below=2pt] at (2.55,2.05) {$A_1$};
  \draw[very thick,Verde] (10.65,2.70) ellipse (0.78 and 0.42);
  \node[below=5pt] at (10.65,2.15) {$A_h$ oculto};
  \draw[flow,draw=Azul] (4.40,1.10) .. controls (3.55,1.45) and (3.0,1.90) .. (2.75,2.30);
  \draw[flow,draw=Verde] (12.75,4.25) .. controls (11.95,3.75) and (11.35,3.30) .. (10.95,2.98);

  \fill[GrisOscuro] (1.50,4.00) circle (2pt);
  \node[left=2pt] at (1.50,4.00) {$e$};
  \draw[GrisOscuro,dashed,thick] (1.50,4.00) circle (0.72);
  \node[above] at (1.50,4.72) {$U_e$};
  \node[draw=GrisOscuro,rounded corners,fill=white,inner sep=3pt]
    at (3.28,3.95) {$U_e\cap\mathcal B(A_h)=\varnothing$};
\end{tikzpicture}%
}

\smallskip
\begin{minipage}{0.92\linewidth}
\small\textbf{Lectura pedagógica.}
La ocultedad se expresa por vecindades de los equilibrios que no intersectan la
cuenca del atractor oculto. La frontera de una cuenca es invariante bajo un
flujo completo, pero identificarla globalmente con $W^s(S)$ exige excluir
otros conjuntos invariantes; el trazo discontinuo representa sólo una
componente candidata.
\end{minipage}
\end{minipage}
\end{center}


En muchos ejemplos una parte de la frontera coincide con la variedad estable
de una silla o de un ciclo inestable. Ésta es una posibilidad estructural, no
una identidad automática. Para demostrar
\(\partial\mathcal B(A)=W^s(S)\) hay que excluir otros conjuntos invariantes en
la frontera y controlar su extensión global.

En el PLL, el ciclo corredor inestable localizado por el mapa de retorno es un
candidato natural a separatriz entre el foco bloqueado y el ciclo corredor
estable. El multiplicador y la figura establecen estabilidad transversal local
y evidencia geométrica; la igualdad global con toda la frontera de cuenca
requiere todavía un argumento en el cilindro.

\subsection{Regiones atrapantes}

\begin{proposition}[Un bloque de entrada produce un conjunto invariante]
Sea \(N\) compacto y supóngase que el campo apunta estrictamente hacia su
interior en \(\partial N\). Para un flujo completo, los conjuntos
\(\varphi_t(N)\), \(t\ge0\), son compactos y anidados, y
\begin{equation}
 A_N=\bigcap_{t\ge0}\varphi_t(N)
 \label{geo:eq:maximal-invariant-trapping}
\end{equation}
es compacto, no vacío, invariante y atrae \(N\).
\end{proposition}

\begin{proof}
El criterio de entrada hace \(N\) positivamente invariante. Por tanto,
\[
 \varphi_{t+s}(N)=\varphi_t(\varphi_s(N))\subset\varphi_t(N).
\]
La
intersección anidada de compactos no vacíos es compacta y no vacía. La
propiedad de flujo da invariancia. Si la distancia de
\(\varphi_t(N)\) a \(A_N\) no tendiera a cero, una sucesión de puntos a
distancia uniforme positiva tendría una subsucesión límite perteneciente a
todas las colas, contradicción.
\end{proof}

Construir un \(N\) así alrededor de un candidato fortalecería sustancialmente
el reporte: reemplazaría la mera permanencia observada por un argumento G3 de
acotación y existencia de un conjunto invariante. Aun así, habría que separar
periodicidad, caos y ocultedad mediante argumentos adicionales.

\section{Métodos topológicos: índice, grado y disparo}
\label{sec:geo-degree}

Esta sección forma el puente más directo con los métodos topológicos de
existencia. El principio central es que un número entero puede permanecer
constante durante una deformación aun cuando no se conozcan explícitamente
las soluciones \cite{Amster2021MetodosTopologicos}.

\subsection{Grado de Brouwer}

\begin{definition}[Grado de Brouwer, datos admisibles]
Sean \(\Omega\subset\R^n\) un abierto acotado,
\(F:\overline\Omega\to\R^n\) continua y
\(y\notin F(\partial\Omega)\). El grado
\(\deg(F,\Omega,y)\in\mathbb Z\) es el entero caracterizado por normalización,
aditividad y invariancia por homotopías que no llevan la frontera sobre \(y\).
\end{definition}

No se reconstruye aquí toda su teoría, pero se usan dos consecuencias.

\begin{proposition}[Existencia e invariancia homotópica]
Con los datos anteriores:
\begin{enumerate}
  \item si \(\deg(F,\Omega,y)\ne0\), existe \(x\in\Omega\) tal que
        \(F(x)=y\);
  \item si \(H:[0,1]\times\overline\Omega\to\R^n\) es continua y
        \(y\notin H(\lambda,\partial\Omega)\) para todo \(\lambda\), entonces
        \(\deg(H(0,\cdot),\Omega,y)=
        \deg(H(1,\cdot),\Omega,y)\).
\end{enumerate}
\end{proposition}

\begin{proof}[Idea]
La primera propiedad se demuestra por contraposición: si \(y\) no pertenece a
la imagen interior, el mapa puede deformarse dentro de
\(\R^n\setminus\{y\}\) hasta uno sin contribución local y el grado es cero. La
segunda es precisamente la invariancia que caracteriza al grado, y la
condición de frontera impide que una raíz entre o salga de \(\Omega\) durante
la deformación.
\end{proof}

Si \(F\in C^1\), todas las raíces \(x_j\) son no degeneradas y finitas,
entonces
\begin{equation}
 \deg(F,\Omega,0)=
 \sum_{F(x_j)=0}\operatorname{sgn}\det DF(x_j).
 \label{geo:eq:degree-sum}
\end{equation}
El sumando es el índice local de la raíz.

\begin{center}
\resizebox{0.96\linewidth}{!}{%
\begin{tikzpicture}[
  >=Latex,
  every node/.style={font=\small},
  map/.style={-{Latex[length=2.5mm]},very thick,draw=GrisOscuro},
  rootp/.style={circle,draw=Azul,fill=AzulClaro,very thick,minimum size=7mm},
  rootm/.style={rectangle,draw=Rojo,fill=NaranjaClaro,very thick,minimum size=7mm}
]
  \node[font=\bfseries] at (3.35,5.10) {Ceros dentro de \(\Omega\)};
  \path[fill=Gris,draw=GrisOscuro,very thick]
    plot[smooth cycle,tension=0.85] coordinates {
      (0.75,1.15) (0.70,3.75) (2.40,4.60)
      (5.30,4.20) (6.05,2.35) (4.95,0.65) (2.10,0.55)
    };
  \node[anchor=north west] at (1.00,4.25) {$\Omega$};
  \node[rootp] (xp) at (2.05,2.80) {$+1$};
  \node[rootm] (xm) at (3.55,1.55) {$-1$};
  \node[rootp] (xpp) at (4.55,3.25) {$+1$};
  \node[below=2pt] at (xp.south) {$x_1$};
  \node[below=2pt] at (xm.south) {$x_2$};
  \node[below=2pt] at (xpp.south) {$x_3$};
  \draw[-{Latex[length=1.8mm]},Azul,thick] (0.88,2.05) -- ++(-0.48,0.35);
  \draw[-{Latex[length=1.8mm]},Azul,thick] (2.95,4.45) -- ++(-0.24,0.42);
  \draw[-{Latex[length=1.8mm]},Azul,thick] (5.70,3.25) -- ++(0.52,0.22);
  \draw[map] (6.55,2.80) -- node[above,align=center]
    {$\widehat F=F/\|F\|$\\sobre \(\partial\Omega\)} (8.25,2.80);

  \node[font=\bfseries] at (10.85,5.10) {Número de vueltas en \(S^1\)};
  \draw[very thick,Azul] (10.85,2.75) circle (1.55);
  \draw[-{Latex[length=2.4mm]},very thick,Azul]
    (12.40,2.75) arc[start angle=0,end angle=310,radius=1.55];
  \draw[GrisOscuro,thin] (8.95,2.75) -- (12.75,2.75);
  \draw[GrisOscuro,thin] (10.85,0.85) -- (10.85,4.65);
  \fill[GrisOscuro] (12.40,2.75) circle (2pt);
  \node[right] at (12.40,2.75) {$1$};
  \node[align=center] at (10.85,2.75)
    {$\deg(F,\Omega,0)$\\$=+1$};

  \node[draw=GrisOscuro,rounded corners,fill=white,inner sep=4pt]
    at (7.10,-0.05)
    {$\displaystyle F(x_j)=0,\qquad
      \deg(F,\Omega,0)
      =\sum_{F(x_j)=0}\operatorname{sgn}\det DF(x_j)
      =(+1)+(-1)+(+1)=1$};
\end{tikzpicture}%
}

\smallskip
\begin{minipage}{0.92\linewidth}
\small\textbf{Lectura pedagógica.}
En dimensión dos, el grado puede verse como el número de vueltas del campo
normalizado sobre la frontera. Los signos locales se suman y certifican la
existencia de al menos un cero cuando el total no es nulo; no determinan por sí
solos estabilidad, cuenca ni caos. El patrón mostrado es esquemático, no un
retrato de fase.
\end{minipage}
\end{center}


\subsection{Ejemplo resuelto: índices de los equilibrios de Chua}

En una región de pendiente \(m\), el reporte verificó
\begin{equation}
 \det J_m=-\alpha[\beta+(\beta+\gamma)m].
 \label{geo:eq:chua-index-det}
\end{equation}
Con los parámetros de~\eqref{int:eq:chua}, la región central \(m=m_0\)
produce \(\det J_{m_0}=-84.035507517056\), mientras las regiones exteriores
\(m=m_1\) producen \(\det J_{m_1}=15.037737580544\). Como los equilibrios no
están sobre \(x=\pm1\), sus índices locales son
\begin{equation}
 \operatorname{ind}(E_0)=-1,\qquad
 \operatorname{ind}(E_+)=\operatorname{ind}(E_-)=+1.
 \label{geo:eq:chua-equilibrium-indices}
\end{equation}
Por tanto, cualquier dominio \(\Omega\) que contenga exactamente esos tres
ceros y satisfaga \(0\notin f(\partial\Omega)\) tiene
\begin{equation}
 \deg(f,\Omega,0)=1.
 \label{geo:eq:chua-degree}
\end{equation}
La conclusión certifica existencia algebraico--topológica de al menos un
equilibrio bajo deformaciones admisibles. No detecta el atractor caótico ni
su cuenca.

El uso del grado sigue siendo válido para el campo Chua continuo aunque no sea
diferenciable en las superficies de conmutación. La fórmula local con
determinantes se aplica porque cada cero está dentro de una región suave.

\subsection{Ejemplo resuelto: índices en MAVPD, Kalman--Fitts y PLL}

Para el Jacobiano MAVPD de~\eqref{int:eq:mavpd-jacobian}, expansión por la
tercera fila da
\begin{equation}
 \det J(y)=\delta\rho(\gamma-3y_1^2).
 \label{geo:eq:mavpd-jacobian-det}
\end{equation}
De aquí,
\begin{equation}
 \operatorname{ind}(E_0)=+1,\qquad
 \operatorname{ind}(E_\pm)=-1,\qquad
 \sum\operatorname{ind}=-1.
 \label{geo:eq:mavpd-indices}
\end{equation}
El signo sólo es índice; no indica por sí mismo estabilidad. Una matriz de
dimensión tres puede tener determinante negativo y, sin embargo, diferentes
combinaciones de dimensiones estable e inestable.

Para Kalman--Fitts, el cálculo simbólico del reporte obtuvo
\(\det Df(X)=0.98191881>0\). El origen, único equilibrio, tiene índice \(+1\)
y el Jacobiano nunca se hace singular. Esto explica a la vez por qué la criba
de puntos perpetuos no produce raíces y por qué el grado del equilibrio no
localiza el ciclo candidato ni establece su ocultedad.

En el PLL, el Jacobiano del campo~\eqref{int:eq:pll} satisface
\begin{equation}
 \det J(x,\theta)=\frac{L}{2T}\cos\theta.
 \label{geo:eq:pll-jacobian-det}
\end{equation}
El equilibrio de fase \(\theta_f\), con \(\cos\theta_f>0\), tiene índice
\(+1\); el equilibrio silla \(\theta_s=\pi-\theta_f\) tiene índice \(-1\).
La suma nula en una celda angular es compatible con la topología cilíndrica,
pero invocar un teorema global de Poincaré--Hopf exigiría además compactificar
o controlar las fronteras en \(x\).

\subsection{Disparo y puntos fijos de Poincaré}

Una órbita periódica puede buscarse como un cero. Si \(P:\Sigma\to\Sigma\) es
un mapa de retorno, se define
\begin{equation}
 F(u)=P(u)-u.
 \label{geo:eq:fixed-point-residual}
\end{equation}
Un grado no nulo de \(F\) en \(U\subset\Sigma\) prueba que \(P\) tiene un
punto fijo dentro de \(U\). Esto convierte el método de disparo en un problema
topológico de existencia.

Para la homotopía Kalman--Fitts
\(\psi_\lambda=(1-\lambda)\operatorname{sign}+\lambda\tanh\), sea
\(P_\lambda\) el mapa de retorno correspondiente. La continuación numérica en
51 niveles es G1. Se convertiría en un argumento de grado si se construyera un
dominio común \(U\) tal que
\begin{equation}
 P_\lambda(u)\ne u
 \quad\text{para todo }u\in\partial U,
 \quad 0\le\lambda\le1,
 \label{geo:eq:kalman-degree-boundary}
\end{equation}
y se verificara
\(\deg(P_0-I,U,0)\ne0\). La invariancia homotópica obligaría a conservar al
menos un punto fijo hasta \(\lambda=1\). El cálculo numérico actual rastrea una
rama; la condición~\eqref{geo:eq:kalman-degree-boundary} excluiría además la
pérdida de todas las soluciones por la frontera.

En un enfoque de dimensión infinita, una solución periódica se formula como
punto fijo de un operador integral sobre un espacio de funciones. El grado de
Leray--Schauder extiende esta estrategia cuando el operador es una
perturbación compacta de la identidad y se tienen cotas a priori. Es una ruta
natural para una fase posterior, pero no debe invocarse sin probar compacidad
y ausencia de soluciones sobre la frontera funcional.

\section{Conjugación, equivalencia orbital y propiedades invariantes}
\label{sec:geo-equivalence}

\begin{definition}[Conjugación topológica]
Dos flujos \(\varphi\) en \(M\) y \(\psi\) en \(N\) son topológicamente
conjugados si existe un homeomorfismo \(h:M\to N\) tal que
\begin{equation}
 h\circ\varphi_t=\psi_t\circ h
 \quad\text{para todo }t.
 \label{geo:eq:topological-conjugacy}
\end{equation}
La igualdad conserva la parametrización temporal.
\end{definition}

\begin{definition}[Equivalencia orbital orientada]
Los flujos son orbitalmente equivalentes si \(h\) lleva cada órbita sobre una
órbita preservando su orientación, pero permite un cambio de reloj
\(\tau(t,x)\), continuo y estrictamente creciente en \(t\):
\begin{equation}
 h(\varphi_t(x))=\psi_{\tau(t,x)}(h(x)).
 \label{geo:eq:orbital-equivalence}
\end{equation}
\end{definition}

\begin{proposition}[Qué transporta una conjugación]
Una conjugación topológica lleva equilibrios a equilibrios, órbitas periódicas
a órbitas periódicas con el mismo periodo, conjuntos invariantes a conjuntos
invariantes, conjuntos \(\omega\)-límite a conjuntos \(\omega\)-límite y
cuencas a cuencas. En particular, la ocultedad definida por
\eqref{geo:eq:hidden-topological} se conserva.
\end{proposition}

\begin{proof}
Cada afirmación se obtiene aplicando \(h\) a la identidad que la define. Por
ejemplo,
\(h(\varphi_T(x))=\psi_T(h(x))=h(x)\) conserva periodicidad; continuidad de
\(h\) y \(h^{-1}\) transporta límites y vecindades; y por tanto
\(h(\mathcal B(A))=\mathcal B(h(A))\).
\end{proof}

Una equivalencia orbital conserva la curva geométrica y sus sentidos, pero no
el periodo ni velocidades. Multiplicar un campo por una función positiva
\begin{equation}
 g(x)=a(x)f(x),\qquad a(x)>0,
 \label{geo:eq:positive-time-rescaling}
\end{equation}
mantiene las órbitas orientadas y cambia el reloj. Por ello, periodos,
velocidades, exponentes de Lyapunov por unidad de tiempo y puntos perpetuos no
son invariantes puramente orbitales.

\subsection{Dos contraejemplos instructivos}

\paragraph{Las tasas lineales no son invariantes topológicas.}
Los sistemas \(\dot x=-x\) y \(\dot y=-2y\) son conjugados mediante
\(h(x)=x|x|\), pues
\[
 h(e^{-t}x)=e^{-2t}h(x).
\]
Ambos tienen la misma organización topológica ---un sumidero y dos
semiórbitas no estacionarias---, aunque sus autovalores son distintos.

\paragraph{Una bola numérica no es un objeto topológico rígido.}
La transformación \(x=Tz+c\) preserva la ocultedad cualitativa, pero la bola
euclidiana se convierte en el elipsoide~\eqref{geo:eq:ball-to-ellipsoid}. Si se
usa otra bola con el mismo radio en \(z\), se ha cambiado el experimento. Esta
distinción explica por qué una equivalencia algebraica de realizaciones no
autoriza a copiar literalmente un sondeo de cuenca.

La sección de puntos perpetuos ofrece otro contraejemplo: su condición es
covariante bajo cambios afines, pero una transformación no lineal introduce
términos Hessianos. En consecuencia, un punto perpetuo es un localizador
dependiente de la representación, no una clase topológica del flujo.

\section{Geometría lineal local y variedades invariantes}
\label{sec:geo-stable-manifolds}

Los resultados de esta sección requieren declarar regularidad e
hiperbolicidad; una gráfica de trayectorias no reemplaza esas hipótesis
\cite{JaurezEtAl2021TGEDI,JaurezEtAl2021TGEDII,HirschSmaleDevaney2013}.

\subsection{Primera variación y flujo tangente}

Si \(f\in C^1\), la dependencia diferenciable respecto de \(x_0\) conduce a
la ecuación variacional
\begin{equation}
 \dot\xi=Df(\varphi_t(x_0))\xi,
 \qquad \Phi(0)=I,\qquad
 \Phi(t)=D\varphi_t(x_0).
 \label{geo:eq:variational}
\end{equation}
Así, \(D\varphi_t\) transporta vectores tangentes. Los exponentes de Lyapunov
y los multiplicadores de Floquet se construyen a partir de esta dinámica en
el fibrado tangente, no a partir de la apariencia de una proyección.

\begin{definition}[Exponente de Lyapunov direccional]
Para \(v\in T_{x_0}M\setminus\{0\}\), se define
\begin{equation}
 \lambda(x_0,v)=\limsup_{t\to\infty}
 \frac{1}{t}\log\frac{\|D\varphi_t(x_0)v\|}{\|v\|}.
 \label{geo:eq:lyapunov-definition}
\end{equation}
El exponente máximo toma el supremo sobre direcciones. En una corrida finita
se estima la cantidad a tiempo \(T\), con renormalizaciones y tolerancias
declaradas; ese estimador no es por sí solo el límite de
\cref{geo:eq:lyapunov-definition} ni una prueba topológica de caos
\cite{Skokos2010,DatserisParlitz2022Nonlinear}.
\end{definition}

\begin{definition}[Equilibrio hiperbólico]
Un equilibrio \(e\) es hiperbólico si ningún autovalor de \(Df(e)\) tiene
parte real cero. Los subespacios espectrales estable e inestable son
\begin{align}
 E^s&=\bigoplus_{\Re\lambda<0}E_\lambda,
 &E^u&=\bigoplus_{\Re\lambda>0}E_\lambda.
 \label{geo:eq:spectral-splitting}
\end{align}
\end{definition}

\begin{proposition}[Hartman--Grobman, alcance local]
Sea \(f\in C^1\) y \(e\) un equilibrio hiperbólico. En una vecindad de \(e\),
el flujo no lineal es topológicamente conjugado al flujo lineal
\(\dot\xi=Df(e)\xi\), restringiendo tiempo y vecindades cuando sea necesario.
\end{proposition}

La demostración completa requiere construir una corrección continua a la
dinámica lineal; la idea clave es que la dicotomía exponencial de la matriz
hiperbólica permite resolver una ecuación integral por punto fijo. El teorema
clasifica la geometría local de órbitas, pero no afirma conjugación
diferenciable, no conserva autovalores y no describe la cuenca global.

\begin{definition}[Conjuntos estable e inestable]
Para un conjunto compacto invariante \(A\),
\begin{align}
 W^s(A)&=\{x:\operatorname{dist}(\varphi_t(x),A)\to0
                    \text{ cuando }t\to+\infty\},\nonumber\\
 W^u(A)&=\{x:\operatorname{dist}(\varphi_t(x),A)\to0
                    \text{ cuando }t\to-\infty\}.
 \label{geo:eq:stable-unstable-sets}
\end{align}
\end{definition}

Para un atractor \(A\), \(W^s(A)\) coincide con su cuenca bajo la convención
de convergencia usada. Para una silla, \(W^s\) suele tener interior vacío y
actuar como separatriz.

\begin{proposition}[Teorema local de la variedad estable]
Sea \(f\in C^r\), \(r\ge1\), y sea \(e\) hiperbólico. Existen variedades
locales \(W^s_{\mathrm{loc}}(e)\) y \(W^u_{\mathrm{loc}}(e)\), de clase
\(C^r\), tangentes en \(e\) a \(E^s\) y \(E^u\), respectivamente. Sus
dimensiones son \(\dim E^s\) y \(\dim E^u\).
\end{proposition}

\begin{proof}[Idea]
En coordenadas espectrales se escribe la ecuación como su parte lineal más un
resto de orden superior. El método de la transformada de gráficas busca la
variedad como gráfica de una función sobre \(E^s\) o \(E^u\). La contracción
producida por la separación espectral da una gráfica invariante; su tangente
en el origen es el subespacio espectral correspondiente.
\end{proof}

\begin{center}
\resizebox{0.96\linewidth}{!}{%
\begin{tikzpicture}[
  >=Latex,
  every node/.style={font=\small},
  manifold/.style={very thick},
  tangent/.style={thick,densely dashed,draw=GrisOscuro},
  phasearrow/.style={-{Latex[length=2.2mm]},very thick}
]
  \draw[rounded corners=6pt,thick,GrisOscuro] (0.15,0.15) rectangle (8.25,5.65);
  \node[font=\bfseries] at (4.20,5.30) {Geometría local de una silla hiperbólica};

  \coordinate (e) at (4.05,2.78);
  \draw[tangent] (1.35,0.30) -- (6.75,5.05);
  \draw[tangent] (0.85,4.20) -- (7.40,1.05);

  \draw[manifold,Azul]
    (1.45,0.30) .. controls (2.25,1.10) and (3.35,2.15) .. (e)
    .. controls (4.75,3.48) and (5.70,4.35) .. (6.60,4.90);
  \draw[manifold,Rojo]
    (0.75,4.42) .. controls (2.10,3.78) and (3.18,3.18) .. (e)
    .. controls (4.95,2.38) and (6.05,1.62) .. (7.48,0.85);

  % Sentido del tiempo sobre las ramas estables.
  \draw[phasearrow,Azul] (2.15,0.95) -- (2.88,1.66);
  \draw[phasearrow,Azul] (5.88,4.55) -- (5.13,3.85);
  % Sentido del tiempo sobre las ramas inestables.
  \draw[phasearrow,Rojo] (3.15,3.18) -- (2.28,3.68);
  \draw[phasearrow,Rojo] (4.98,2.35) -- (5.90,1.78);

  \fill[GrisOscuro] (e) circle (2.6pt);
  \node[below right] at (e) {$e$};
  \node[Azul,fill=white,inner sep=2pt] at (2.15,1.90) {$W^s_{\mathrm{loc}}(e)$};
  \node[Rojo,fill=white,inner sep=2pt] at (6.35,2.35) {$W^u_{\mathrm{loc}}(e)$};
  \node[Azul,fill=white,inner sep=2pt] at (6.22,4.52) {$E^s$};
  \node[Rojo,fill=white,inner sep=2pt] at (6.75,1.30) {$E^u$};

  \draw[rounded corners=6pt,thick,GrisOscuro] (8.65,0.15) rectangle (13.95,5.65);
  \node[font=\bfseries] at (11.30,5.30) {Definición dinámica};
  \node[align=left,text width=4.55cm] at (11.30,4.08) {
    \(x\in W^s(e)\) si\\[-1mm]
    \(\varphi_t(x)\to e\) cuando \(t\to+\infty\).};
  \node[align=left,text width=4.55cm] at (11.30,2.78) {
    \(x\in W^u(e)\) si\\[-1mm]
    \(\varphi_t(x)\to e\) cuando \(t\to-\infty\).};
  \draw[GrisOscuro,dashed] (9.10,2.00) -- (13.50,2.00);
  \node[align=center] at (11.30,1.30)
    {$T_eW^s_{\mathrm{loc}}=E^s$,\qquad
     $T_eW^u_{\mathrm{loc}}=E^u$};
  \node[align=center,text=GrisOscuro,font=\scriptsize,text width=4.2cm]
    at (11.30,0.70)
    {tangencia local\\\(\neq\) descripción global};
\end{tikzpicture}%
}

\smallskip
\begin{minipage}{0.92\linewidth}
\small\textbf{Lectura pedagógica.}
Los subespacios espectrales describen las tangentes en el equilibrio; las
variedades invariantes son objetos no lineales. Las flechas azules llegan a la
silla en tiempo futuro y las rojas se alejan de ella. El dibujo es un corte
local bidimensional y no fija la forma global de las separatrices.
\end{minipage}
\end{center}


\subsection{Variedad central y el límite de la linealización}

Si \(Df(e)\) tiene autovalores sobre el eje imaginario, el equilibrio no es
hiperbólico y Hartman--Grobman no decide su dinámica. Se añade entonces el
subespacio
\begin{equation}
 E^c=\bigoplus_{\Re\lambda=0}E_\lambda,\qquad
 T_eM=E^s\oplus E^c\oplus E^u.
 \label{geo:eq:center-splitting}
\end{equation}

\begin{proposition}[Teorema local de la variedad central, alcance]
Sea \(f\in C^r\), \(r\ge1\), con \(f(e)=0\). Existe localmente una variedad
invariante \(W^c_{\mathrm{loc}}(e)\), tangente a \(E^c\) en \(e\), cuya
regularidad está limitada por la de \(f\). La dinámica local relevante que no
se contrae ni expande exponencialmente puede reducirse a ella. En general la
variedad central no es única fuera del orden local necesario.
\end{proposition}

La construcción se realiza, como en la variedad estable, escribiendo una
gráfica y resolviendo una ecuación de invariancia; la ausencia de contracción
en \(E^c\) explica por qué la unicidad y la regularidad son más delicadas.
Conocer sólo \(Df(e)\) no basta. Por ejemplo,
\begin{equation}
 \dot x=-x^3
 \qquad\text{y}\qquad
 \dot x=x^3
 \label{geo:eq:center-counterexample}
\end{equation}
tienen la misma linealización nula en el origen, pero el primero es estable y
el segundo inestable.

\subsection{Bifurcaciones locales elementales}

En una familia escalar \(\dot x=f(x,\mu)\), una silla--nodo genérica en
\((x_*,\mu_*)\) requiere
\begin{equation}
 f=0,\qquad f_x=0,\qquad f_{xx}\ne0,\qquad f_\mu\ne0.
 \label{geo:eq:saddle-node-conditions}
\end{equation}
La forma normal \(\dot x=\mu-x^2\) permite comprobar que dos equilibrios
colisionan y desaparecen. Una bifurcación de Hopf exige en cambio un par simple
de autovalores imaginarios, cruce transversal y no degeneración del primer
coeficiente de Lyapunov. La igualdad espectral localiza sólo un candidato al
evento; las condiciones restantes deciden si nace una órbita y su estabilidad
\cite{Kuznetsov2023Bifurcation,JaurezEtAl2021TGEDII}.

La igualdad de Routh--Hurwitz usada para la frontera candidata de Hopf MAVPD
coloca un par sobre el eje imaginario y sugiere una variedad central
bidimensional. Para afirmar una bifurcación de Hopf deben comprobarse además
la simplicidad del par, la transversalidad del cruce y una condición de no
degeneración, normalmente mediante el primer coeficiente de Lyapunov. El
reporte no presenta aún esas tres pruebas, por lo que conserva correctamente
la expresión «frontera candidata».

\subsection{Alcance caso por caso}

\begin{table}[H]
\centering
\caption{Uso correcto del teorema de la variedad estable en los expedientes.}
\label{geo:tab:stable-manifold-scope}
\small
\begin{tabularx}{\textwidth}{@{}L{0.17\textwidth}YY@{}}
\toprule
Caso & Aplicación disponible & Límite que debe declararse \\
\midrule
Chua PWL & En cada equilibrio interior a una región suave, si su espectro es
hiperbólico, existen variedades locales clásicas. & Al cruzar \(x=\pm1\), la
continuación global es por tramos; en un roce no puede suponerse suavidad del
mapa de paso. \\
MAVPD & El campo es suave; cada equilibrio hiperbólico admite variedades con
dimensiones dadas por su espectro. & El determinante o Routh--Hurwitz aislado
no determina toda la geometría global ni una conexión homoclínica. \\
Kalman--Fitts & Si el espectro archivado del origen queda estrictamente en el
semiplano izquierdo, su variedad estable local tiene dimensión cuatro. Un
ciclo hiperbólico posee variedades en torno de su órbita. & El
origen no tiene rama inestable que conduzca al ciclo candidato. Los 48 sondeos
disponibles no identifican la frontera global de su cuenca; ésta debe buscarse
en otros conjuntos invariantes. \\
PLL & El foco y la silla se estudian en cartas del cilindro. La silla posee
separatrices unidimensionales. & Una curva dibujada en el recubrimiento debe
pegarse módulo \(2\pi\); no toda separatriz local es toda la frontera global.
\\
\bottomrule
\end{tabularx}
\end{table}

\subsection{Divergencia y cambio de volumen}

La geometría diferencial permite seguir elementos de volumen mediante el
determinante de \(D\varphi_t\).

\begin{proposition}[Fórmula de Liouville]
Sea \(f\in C^1\) y supóngase definido el flujo. Entonces
\begin{equation}
 \det D\varphi_t(x)=
 \exp\!\left(
   \int_0^t\operatorname{div}f(\varphi_s(x))\,ds
 \right).
 \label{geo:eq:liouville-volume}
\end{equation}
\end{proposition}

\begin{proof}
La matriz \(\Phi=D\varphi_t(x)\) satisface~\eqref{geo:eq:variational}. La
fórmula de Jacobi da
\[
 \frac{d}{dt}\det\Phi=\operatorname{tr}(Df)\det\Phi.
\]
Como
\(\det\Phi(0)=1\) y \(\operatorname{tr}(Df)=\operatorname{div}f\), integrar
la ecuación escalar produce~\eqref{geo:eq:liouville-volume}.
\end{proof}

Para MAVPD,
\begin{equation}
 \operatorname{div}f(y)=\delta(\gamma-3y_1^2)-\xi.
 \label{geo:eq:mavpd-divergence}
\end{equation}
Su signo puede variar con \(y_1\), por lo que una afirmación de disipatividad
uniforme necesita acotar la región visitada. Incluso
\(\operatorname{div}f<0\) en todas partes sólo implica contracción de volumen;
el sistema \(\dot x=-x,\dot y=-y\) contrae área y no es caótico. Contracción,
acotación y caos son tres afirmaciones diferentes.

\subsection{Formas diferenciales, flujo y teorema de Stokes}

En una variedad orientada de dimensión \(n\), una forma diferencial de grado
\(k\) asigna alternadamente un número a \(k\) vectores tangentes. En
\(\R^n\), la forma de volumen estándar es
\begin{equation}
 \Omega=dx_1\wedge\cdots\wedge dx_n.
 \label{geo:eq:volume-form}
\end{equation}
La contracción interior \(\iota_f\Omega\) es una forma de grado \(n-1\) que
mide flujo a través de hipersuperficies. La derivada de Lie de \(\Omega\) a lo
largo de \(f\) satisface
\begin{equation}
 \mathcal L_f\Omega=(\operatorname{div}f)\Omega,\qquad
 d(\iota_f\Omega)=(\operatorname{div}f)\Omega.
 \label{geo:eq:cartan-divergence}
\end{equation}
La segunda igualdad usa \(d\Omega=0\) y la fórmula de Cartan.

\begin{proposition}[Stokes como balance de flujo]
Sea \(N\) una región compacta orientada con frontera suave y \(f\in C^1\).
Entonces
\begin{equation}
 \int_{\partial N}\iota_f\Omega
 =\int_N d(\iota_f\Omega)
 =\int_N(\operatorname{div}f)\Omega.
 \label{geo:eq:stokes-flux}
\end{equation}
\end{proposition}

\begin{proof}
La primera igualdad es el teorema de Stokes aplicado a la forma
\(\iota_f\Omega\); la segunda es~\eqref{geo:eq:cartan-divergence}.
\end{proof}

Este lenguaje enlaza el criterio puntual de entrada
\eqref{geo:eq:strict-inward} con un balance integral. Sin embargo, flujo total
negativo a través de \(\partial N\) no implica que el campo apunte hacia
dentro en cada punto: regiones de salida pueden compensarse con regiones de
entrada. Por ello, Stokes no sustituye la desigualdad puntual necesaria para
una región atrapante.

En el PLL se usa la forma global \(dx\wedge d\theta\) sobre el cilindro; en
Chua, el campo es continuo pero sólo suave por regiones, de modo que el balance
se aplica por piezas y las contribuciones internas deben emparejarse en
\(x=\pm1\). Esta formulación ofrece una conexión directa entre el libro de
topología y geometría diferencial y los cálculos de divergencia del reporte.

\section{Mapas de Poincaré, monodromía y teoría de Floquet}
\label{sec:geo-poincare-floquet}

El paso flujo--retorno reduce la estabilidad orbital a dinámica discreta, pero
sólo después de verificar transversalidad, existencia del primer retorno y
dependencia regular respecto del punto inicial
\cite{JaurezEtAl2021TGEDII,HirschSmaleDevaney2013,Kuznetsov2023Bifurcation}.

\subsection{Reducir una órbita periódica a un punto fijo}

Sea \(\gamma\) una órbita periódica de periodo mínimo \(T\) y
\(p\in\gamma\). Una sección local \(\Sigma\) de codimensión uno es transversal
si
\begin{equation}
 T_pM=T_p\Sigma\oplus\operatorname{span}\{f(p)\}.
 \label{geo:eq:transverse-section}
\end{equation}
Bajo \(f\in C^1\) y transversalidad, el teorema de la función implícita
produce un tiempo de primer retorno \(\tau(x)\) para \(x\) cerca de \(p\), y
el mapa
\begin{equation}
 P(x)=\varphi_{\tau(x)}(x),\qquad P:\Sigma\to\Sigma.
 \label{geo:eq:poincare-map}
\end{equation}
La órbita \(\gamma\) se convierte en el punto fijo \(P(p)=p\).

\begin{proposition}[Estabilidad mediante el mapa de retorno]
Si todos los autovalores de \(DP(p)\) tienen módulo menor que uno, \(p\) es un
punto fijo localmente asintóticamente estable de \(P\) y \(\gamma\) es una
órbita periódica orbitalmente asintóticamente estable. Si alguno tiene módulo
mayor que uno, la órbita es inestable.
\end{proposition}

\begin{proof}[Idea]
Si el radio espectral es menor que uno, existe una norma adaptada en la cual
\(P\) es contractivo en una vecindad, y sus iterados convergen a \(p\). El
flujo entre retornos transporta esa convergencia alrededor de toda la órbita.
Una dirección propia con módulo mayor que uno se expande en iteraciones del
retorno y produce inestabilidad transversal.
\end{proof}

La sección debe registrar orientación de cruce y eliminar el transitorio.
Contar puntos cercanos generados varias veces por interpolación no equivale a
iterar un mapa de primer retorno.

\subsection{Floquet y los multiplicadores}

Sobre \(\gamma(t)\), la ecuación variacional posee coeficientes periódicos
\begin{equation}
 \dot\xi=A(t)\xi,\qquad A(t)=Df(\gamma(t)),\qquad A(t+T)=A(t).
 \label{geo:eq:periodic-variational}
\end{equation}
La matriz \(M=\Phi(T)\) es la matriz de monodromía y sus autovalores son los
multiplicadores de Floquet. La teoría de Floquet factoriza
\begin{equation}
 \Phi(t)=Q(t)e^{tR},\qquad Q(t+T)=Q(t).
 \label{geo:eq:floquet-factorization}
\end{equation}

\begin{proposition}[Relación Poincaré--Floquet]
Para una órbita periódica regular de un sistema autónomo \(C^1\), \(M\) tiene
un multiplicador trivial \(1\), asociado con la dirección \(f(\gamma(t))\).
Los otros \(n-1\) multiplicadores coinciden, contando multiplicidad, con los
autovalores de \(DP(p)\).
\end{proposition}

\begin{proof}[Idea]
Derivar \(\gamma(t+s)\) respecto de \(s\) en \(s=0\) muestra que
\(f(\gamma(t))\) resuelve la ecuación variacional y vuelve a sí mismo después
de \(T\); de ahí el multiplicador \(1\). El cociente por la dirección tangente
se identifica con el espacio transversal \(T_p\Sigma\), donde actúa
\(DP(p)\).
\end{proof}

\begin{center}
\begin{minipage}{0.96\linewidth}
\centering
\begin{tikzpicture}[
  scale=0.90,
  transform shape,
  >=Latex,
  every node/.style={font=\small},
  flow/.style={-{Latex[length=2.3mm]},very thick,draw=Azul},
  point/.style={circle,fill=GrisOscuro,inner sep=1.9pt}
]
  \draw[rounded corners=6pt,thick,GrisOscuro] (0.10,0.15) rectangle (7.35,5.75);
  \node[font=\bfseries] at (3.72,5.40) {Órbita y sección transversal};
  \draw[very thick,Azul] (2.75,2.78) ellipse (2.15 and 1.45);
  \draw[flow] (2.12,4.16) -- (3.00,4.20);
  \node[Azul] at (1.05,1.18) {$\gamma$};

  \draw[very thick,Rojo] (3.45,2.78) -- (6.65,2.78);
  \node[Rojo,above] at (6.25,2.78) {$\Sigma$};
  \node[point] (p) at (4.90,2.78) {};
  \node[above=2pt] at (p.north) {$p=P(p)$};
  \node[point] (x) at (3.82,2.78) {};
  \node[below=2pt] at (x.south) {$x$};
  \node[point] (px) at (4.38,2.78) {};
  \node[below=2pt] at (px.south) {$P(x)$};
  \draw[flow,draw=GrisOscuro]
    (x) .. controls (6.40,5.25) and (0.25,5.30) .. (0.35,2.75)
        .. controls (0.30,0.20) and (6.20,0.15) .. (px);
  \node[fill=white,inner sep=2pt] at (3.55,0.52)
    {$P(x)=\varphi_{\tau(x)}(x)$};
  \node[align=center,text=GrisOscuro] at (3.72,4.82)
    {$T_pM=T_p\Sigma\oplus\operatorname{span}\{f(p)\}$};

  \draw[rounded corners=6pt,thick,GrisOscuro] (7.75,0.15) rectangle (13.95,5.75);
  \node[font=\bfseries] at (10.85,5.40) {Multiplicadores};
  \coordinate (c) at (10.65,2.82);
  \draw[GrisOscuro,thin] (8.55,2.82) -- (13.15,2.82);
  \draw[GrisOscuro,thin] (10.65,0.72) -- (10.65,4.92);
  \draw[very thick,Azul] (c) circle (1.60);
  \node[below right,text=Azul] at (11.80,1.60) {$|\mu|=1$};

  \fill[GrisOscuro] (12.25,2.82) circle (2.5pt);
  \node[above right] at (12.25,2.82) {$\mu_1=1$ tangencial};
  \fill[Verde] (10.00,3.48) circle (2.5pt);
  \fill[Verde] (11.15,1.92) circle (2.5pt);
  \node[align=right,Verde] at (9.25,4.10) {transversales\\contractivos};
  \draw[Rojo,very thick,fill=white] (12.78,3.78) circle (2.8pt);
  \node[align=left,Rojo] at (12.58,4.52) {alternativa\\expansiva};

  \node[draw=GrisOscuro,rounded corners,fill=white,inner sep=4pt]
    at (10.85,0.45)
    {$\operatorname{spec}\Phi(T)=\{1\}\cup\operatorname{spec}DP(p)$};
\end{tikzpicture}%
\end{minipage}

\smallskip
\begin{minipage}{0.92\linewidth}
\small\textbf{Lectura pedagógica.}
La órbita periódica se reduce al punto fijo $P(p)=p$ de un retorno sobre la
sección. La monodromía siempre conserva la dirección tangente con multiplicador
$1$; los $n-1$ multiplicadores restantes son los autovalores de $DP(p)$.
Los puntos interiores y el punto exterior representan criterios alternativos,
no un mismo espectro medido.
\end{minipage}
\end{center}


Una órbita periódica es hiperbólica si ninguno de sus multiplicadores no
triviales está sobre el círculo unitario. La proximidad numérica de un
multiplicador a \(1\) exige un análisis de error: puede indicar bifurcación,
convergencia lenta o simplemente integración insuficiente.

\subsection{Ejemplo resuelto: multiplicador transversal del PLL}

En dimensión dos hay un solo multiplicador no trivial. Por
\eqref{geo:eq:liouville-volume}, el producto de los dos multiplicadores es
\begin{equation}
 \det\Phi(T)=
 \exp\!\left(\int_0^T\operatorname{div}f(\gamma(t))\,dt\right).
\end{equation}
Como el multiplicador tangente es \(1\), el transversal es exactamente
\begin{equation}
 \mu_\perp=
 \exp\!\left[\int_0^T
 \left(-\frac1{T_{\!f}}
 -\frac{\tau_2L}{2T_{\!f}}\cos\theta(t)\right)dt\right],
 \qquad T_{\!f}=\tau_1+\tau_2.
 \label{geo:eq:pll-floquet-divergence}
\end{equation}
Aquí se usa \(T_f\) para el parámetro del filtro y \(T\) para el periodo de la
órbita. Si la integral es negativa, \(0<\mu_\perp<1\) y el ciclo es estable;
si es positiva, \(\mu_\perp>1\) y es inestable. Esta identidad explica
geométricamente los multiplicadores reportados para el ciclo corredor estable
y la órbita separadora inestable.

\subsection{Lectura correcta de los expedientes}

\begin{table}[H]
\centering
\caption{Qué aporta Poincaré en cada caso y qué queda pendiente.}
\label{geo:tab:poincare-cases}
\small
\begin{tabularx}{\textwidth}{@{}L{0.17\textwidth}YY@{}}
\toprule
Caso & Lectura disponible & Siguiente verificación geométrica \\
\midrule
Chua PWL & La recurrencia y las nubes apoyan un conjunto caótico; las órbitas
cruzan regiones PWL. & Declarar sección y orientación; propagar la primera
variación por tramos; estudiar si existe una órbita periódica silla y sus
variedades. \\
MAVPD \(\xi=3.1\) & Los últimos retornos colapsan numéricamente a un punto y el
periodo es estable: evidencia fuerte de ciclo de periodo uno. & Encerrar el
punto fijo de \(P\) y sus multiplicadores no triviales. \\
MAVPD caótico & La nube de Poincaré, los exponentes y la prueba \(0\)--\(1\)
son diagnósticos G1 compatibles con caos. & Buscar estructura hiperbólica,
órbitas periódicas silla y cruces de variedades; una nube no es una herradura.
\\
Kalman--Fitts & El mapa de conmutación de la aproximación signo tiene retorno
de periodo dos; la homotopía llega a un ciclo suave estable. & Construir una
sección tridimensional para el sistema final y calcular los tres
multiplicadores transversales. \\
PLL & El retorno escalar sobre \(\theta=0\pmod{2\pi}\) localiza ciclos
corredores estable e inestable con multiplicadores. & Encerrar los puntos
fijos y demostrar qué componente de la frontera de cuenca genera el ciclo
inestable. \\
\bottomrule
\end{tabularx}
\end{table}

Para Chua, la falta de diferenciabilidad de \(f\) en \(x=\pm1\) impide citar
sin comentario la versión global \(C^1\) de Floquet. Con cruces transversales
se integra la ecuación variacional en cada región y se tratan explícitamente
los tiempos de evento; en un roce, la diferenciabilidad del retorno puede
fallar.

\subsection{Poincaré--Bendixson y por qué la dimensión importa}

\begin{proposition}[Poincaré--Bendixson, versión planar]
Sea \(f\in C^1\) en un abierto del plano y sea una semiórbita positiva
contenida en un compacto. Si su conjunto \(\omega\)-límite no contiene
equilibrios, entonces ese conjunto \(\omega\)-límite es una órbita periódica.
\end{proposition}

\begin{proof}[Idea geométrica]
Se toma una sección transversal pequeña. Si la órbita regresara desordenada,
la unicidad y el orden lineal de los puntos sobre la sección obligarían a
cruces de trayectorias. Los retornos quedan ordenados y convergen a un punto
fijo del retorno, que genera una órbita periódica. La demostración completa
controla la recurrencia y los arcos de Jordan.
\end{proof}

El teorema explica por qué un flujo autónomo planar suave no alberga un
atractor extraño en una región compacta bajo esas hipótesis. No se aplica a
Chua ni MAVPD, que son tridimensionales, ni a Kalman--Fitts, que es
cuatridimensional. Tampoco se traslada literalmente a toda superficie: el
flujo lineal de pendiente irracional en el toro no tiene equilibrios ni
órbitas periódicas y puede ser denso.

El PLL vive en un cilindro no compacto, no en el plano entero. Si la dinámica
queda confinada a un anillo compacto y se verifican las hipótesis de una
versión anular de Poincaré--Bendixson, su organización por focos, sillas y
ciclos corredores es la esperada. El devanado angular debe conservarse: una
órbita que aumenta \(2\pi\) por vuelta es periódica en el cilindro, aunque no
sea cerrada en el recubrimiento. Esta limitación dimensional hace al PLL un
banco para cuencas y separatrices, no el primer banco para demostrar una
herradura de flujo.

\begin{remark}[Límites del teorema]
Compacidad, unicidad, autonomía, dimensión y topología del espacio son
hipótesis sustantivas. Un mapa discreto bidimensional puede ser caótico; una
ecuación planar forzada se convierte en un sistema autónomo de dimensión tres
al añadir la fase; y un campo no único puede tener cruces que destruyen el
argumento planar.
\end{remark}

\section{Homoclinicidad, herraduras y significados de caos}
\label{sec:geo-homoclinic-chaos}

Esta sección delimita una agenda. No se afirma que los expedientes actuales
ya contengan una prueba homoclínica. Las definiciones y el mecanismo de
herradura se estudian en
\cite{HirschSmaleDevaney2013,DatserisParlitz2022Nonlinear}; su uso en un caso
requiere además cerrar numérica o analíticamente las hipótesis geométricas.

\subsection{Conexiones entre variedades}

\begin{definition}[Órbita homoclínica y heteroclínica]
Sea \(S\) un equilibrio o una órbita periódica hiperbólica. Una órbita no
trivial es homoclínica a \(S\) si está contenida en
\(W^s(S)\cap W^u(S)\). Una conexión de \(S_-\) a \(S_+\) es heteroclínica si
está en \(W^u(S_-)\cap W^s(S_+)\). Una intersección es transversal en \(q\)
si
\begin{equation}
 T_qW^s+T_qW^u=T_qM
 \label{geo:eq:transverse-intersection}
\end{equation}
con la adaptación correspondiente en una sección de Poincaré.
\end{definition}

Una proximidad visual entre dos curvas no prueba que se intersecten; una
intersección numérica no prueba transversalidad; y una conexión no transversal
no activa automáticamente el teorema de herradura.

\begin{proposition}[Teorema de Smale--Birkhoff, versión de agenda]
Sea \(P\) un difeomorfismo \(C^1\) con un punto periódico hiperbólico. Si sus
variedades estable e inestable poseen un punto homoclínico transversal,
entonces algún iterado de \(P\) posee un conjunto compacto invariante sobre el
cual la dinámica contiene una herradura y es conjugada a un desplazamiento
simbólico sobre un número finito de símbolos.
\end{proposition}

La demostración completa construye tiras que se estiran, contraen y cruzan,
define rectángulos de Markov y codifica cada órbita por la sucesión de tiras
visitadas. No se reproduce como «demostración breve» porque precisamente esa
construcción global es lo que tendría que verificarse en un caso del reporte.
El teorema sí fija con claridad las hipótesis que faltan.

\begin{center}
\begin{minipage}{\linewidth}
\centering
\resizebox{0.96\linewidth}{!}{%
\begin{tikzpicture}[
  >=Latex,
  every node/.style={font=\small},
  ws/.style={very thick,draw=Azul},
  wu/.style={very thick,draw=Rojo},
  flow/.style={-{Latex[length=2.3mm]},very thick}
]
  \draw[rounded corners=6pt,thick,GrisOscuro] (0.10,0.10) rectangle (6.35,5.75);
  \node[font=\bfseries] at (3.22,5.42) {Retorno homoclínico en \(\Sigma\)};

  \coordinate (p) at (2.35,2.10);
  \draw[ws]
    (2.18,0.38) .. controls (1.85,1.35) and (2.55,3.25) .. (2.25,5.10);
  \draw[wu]
    (0.45,2.36) .. controls (1.18,2.05) and (1.70,2.02) .. (p)
    .. controls (4.15,1.95) and (5.55,2.55) .. (5.20,3.72)
    .. controls (4.90,4.78) and (3.35,4.65) .. (2.30,4.25)
    .. controls (1.30,3.88) and (0.80,3.32) .. (0.55,2.80);
  \fill[GrisOscuro] (p) circle (2.7pt);
  \node[below right] at (p) {$p$};
  \coordinate (q) at (2.30,4.25);
  \fill[GrisOscuro] (q) circle (2.7pt);
  \node[left=5pt] at (q) {$q$};
  \node[Azul,fill=white,inner sep=2pt] at (1.50,4.72) {$W^s(p)$};
  \node[Rojo,fill=white,inner sep=2pt] at (4.62,4.45) {$W^u(p)$};
  \draw[flow,Rojo] (3.10,2.08) -- (3.82,2.18);
  \draw[flow,Rojo] (4.80,3.28) -- (4.62,3.90);
  \draw[flow,Azul] (2.22,3.25) -- (2.27,2.62);
  \draw[flow,Azul] (2.25,4.82) -- (2.28,4.43);
  \node[align=center,text=GrisOscuro,font=\scriptsize] at (3.25,0.72)
    {$q\in W^s(p)\cap W^u(p)$\\
      $T_qW^s+T_qW^u=T_q\Sigma$};

  \draw[-{Latex[length=2.7mm]},very thick,GrisOscuro]
    (6.70,2.90) -- node[above,align=center] {Smale--\\Birkhoff} (8.05,2.90);

  \draw[rounded corners=6pt,thick,GrisOscuro] (8.35,0.10) rectangle (13.95,5.75);
  \node[font=\bfseries] at (11.15,5.42) {Estirar, contraer y plegar};
  \draw[line width=16pt,Azul!60,rounded corners=12pt]
    (9.25,0.42) -- (9.25,5.00) -- (12.45,5.00) -- (12.45,0.42);
  \draw[line width=7pt,white,rounded corners=9pt]
    (9.25,0.42) -- (9.25,5.00) -- (12.45,5.00) -- (12.45,0.42);
  \draw[very thick,GrisOscuro] (8.78,0.85) rectangle (12.92,4.38);
  \node[anchor=north west] at (8.88,4.28) {$R$};
  \node[font=\bfseries] at (9.25,2.35) {$R_0$};
  \node[font=\bfseries] at (12.45,2.35) {$R_1$};
  \draw[-{Latex[length=2.3mm]},thick,GrisOscuro]
    (10.05,1.35) -- (10.05,0.72);
  \draw[-{Latex[length=2.3mm]},thick,GrisOscuro]
    (11.65,1.35) -- (11.65,0.72);
  \node[draw=GrisOscuro,rounded corners,fill=white,inner sep=3pt]
    at (11.15,0.42)
    {$\cdots010011\cdots\in\{0,1\}^{\mathbb Z}$};
\end{tikzpicture}%
}

\smallskip
\begin{minipage}{0.92\linewidth}
\small\textbf{Lectura pedagógica.}
Para un mapa de retorno, una intersección homoclínica transversal de las
variedades de un punto hiperbólico permite construir tiras que se codifican por
símbolos. La flecha representa una implicación teórica bajo las hipótesis de
Smale--Birkhoff; no afirma que la intersección ya haya sido certificada en Chua
o MAVPD.
\end{minipage}
\end{minipage}
\end{center}


Para un flujo tridimensional como Chua o MAVPD, la búsqueda puede hacerse en
un mapa de Poincaré bidimensional: las curvas estable e inestable de un punto
silla periódico pueden cruzarse transversalmente. En Kalman--Fitts, la
sección es tridimensional y la geometría puede involucrar curvas y
superficies. En un flujo autónomo bidimensional como PLL sobre el cilindro, la
unicidad impide que separatrices distintas se crucen transversalmente como
trayectorias del flujo; su teoría natural es la organización por ciclos,
sillas y fronteras, junto con versiones de Poincaré--Bendixson en una región
compacta apropiada.

\subsection{Dinámica simbólica y entropía}

\begin{definition}[Desplazamiento completo en dos símbolos]
Sea
\begin{equation}
 \Sigma_2=\{0,1\}^{\mathbb Z},\qquad
 (\sigma s)_j=s_{j+1}.
 \label{geo:eq:full-shift}
\end{equation}
Con la topología producto, \(\Sigma_2\) es compacto. Una secuencia codifica qué
rectángulo de una herradura visita una órbita en cada iteración. Decir que
\(P^k|_\Lambda\) es conjugado a \(\sigma\) significa que existe un
homeomorfismo entre \(\Lambda\) y \(\Sigma_2\) que entrelaza exactamente ambos
mapas.
\end{definition}

Las sucesiones periódicas son densas en \(\Sigma_2\), y el desplazamiento es
topológicamente transitivo. Además, su entropía topológica es \(\log2\). Como
la entropía es invariante por conjugación y
\(h_{\mathrm{top}}(P^k)=k\,h_{\mathrm{top}}(P)\), una herradura conjugada en
un iterado proporciona la cota
\begin{equation}
 h_{\mathrm{top}}(P)\ge\frac{\log2}{k},
 \label{geo:eq:horseshoe-entropy-bound}
\end{equation}
porque \(h_{\mathrm{top}}(P^k|_\Lambda)=\log2\). Éste es un puente demostrable
entre geometría de variedades, codificación simbólica y complejidad
topológica; una estimación espectral o una FFT no proporciona esta cota.

\begin{definition}[Transitividad y caos de Devaney]
Un mapa continuo \(P:\Lambda\to\Lambda\) es topológicamente transitivo si para
cualesquiera abiertos relativos no vacíos \(U,V\subset\Lambda\) existe
\(n\ge0\) tal que \(P^n(U)\cap V\ne\varnothing\). Es caótico en el sentido de
Devaney si es transitivo y sus puntos periódicos son densos; usualmente se
incluye sensibilidad a condiciones iniciales, que se deduce de las dos
propiedades anteriores cuando \(\Lambda\) es un espacio métrico infinito sin
puntos aislados.
\end{definition}

La transitividad es una afirmación sobre todos los pares de abiertos, no sobre
una única trayectoria larga. La densidad de órbitas periódicas tampoco se
demuestra encontrando unas cuantas. Por eso la conjugación simbólica es mucho
más fuerte que una nube visual: entrega ambas propiedades sobre un conjunto
invariante bien definido.

\subsection{No hay una sola definición operativa de caos}

\begin{table}[H]
\centering
\caption{Nociones que no deben intercambiarse sin teorema.}
\label{geo:tab:chaos-notions}
\small
\begin{tabularx}{\textwidth}{@{}L{0.20\textwidth}Y Y@{}}
\toprule
Noción & Contenido & Evidencia o prueba necesaria \\
\midrule
Exponente máximo positivo & Expansión asintótica de vectores tangentes para
una órbita o medida, bajo hipótesis de existencia. & Convergencia y robustez
numérica son G1; una prueba requiere controlar el límite y la órbita o medida.
\\
Transitividad topológica & Para abiertos no vacíos \(U,V\) del conjunto
invariante existe \(t>0\) con \(\varphi_t(U)\cap V\ne\varnothing\). & No se
deduce de que una sola órbita parezca llenar una proyección. \\
Caos de Devaney & Transitividad y densidad de órbitas periódicas; en espacios
perfectos, la sensibilidad se sigue bajo hipótesis estándar. & Exige un
conjunto invariante y argumentos sobre todos sus abiertos. \\
Entropía topológica positiva & Crecimiento exponencial de órbitas
distinguibles. & Puede acotarse mediante una herradura o técnicas rigurosas;
una FFT de banda ancha no la calcula. \\
Herradura & Conjunto hiperbólico con codificación simbólica. & Intersección
homoclínica transversal o cobertura de conos/rectángulos verificada. \\
\bottomrule
\end{tabularx}
\end{table}

\subsection{Protocolo propuesto para Chua y MAVPD}

Para convertir los diagnósticos actuales en una investigación geométrica de
caos, el orden correcto es:

\begin{enumerate}
  \item fijar una sección transversal y un mapa de primer retorno bien
        definido en un dominio \(D\subset\Sigma\);
  \item localizar órbitas periódicas de tipo silla como ceros de
        \(P^k(u)-u\), con una condición que excluya periodos divisores;
  \item encerrar el punto periódico y los autovalores de \(DP^k\), evitando el
        círculo unitario;
  \item parametrizar segmentos locales de \(W^s\) y \(W^u\) mediante los
        autovectores y continuarlos con control de error;
  \item encerrar una intersección y verificar
        \eqref{geo:eq:transverse-intersection}, por ejemplo con aritmética de
        intervalos y conos;
  \item sólo entonces invocar Smale--Birkhoff y declarar el subconjunto y el
        iterado para los cuales se ha probado dinámica simbólica.
\end{enumerate}

En Chua deben añadirse eventos exactos en \(x=\pm1\) y excluir roces dentro
del dominio del retorno. En MAVPD, la suavidad evita esa dificultad, por lo
que el candidato caótico local es el primer banco conveniente para desarrollar
el protocolo.

\section{Lectura geométrica unificada de los cuatro casos}
\label{sec:geo-four-cases}

\begin{table}[H]
\centering
\caption{Objetos geométricos ya presentes y preguntas nuevas.}
\label{geo:tab:four-cases-map}
\small
\begin{tabularx}{\textwidth}{@{}L{0.16\textwidth}Y Y@{}}
\toprule
Caso & Objetos ya construidos & Pregunta geométrica prioritaria \\
\midrule
Chua PWL & Tres equilibrios, dos superficies de conmutación, simetría impar,
nubes recurrentes, candidato caótico y sondeos de cuenca. & ¿Qué conjuntos
invariantes forman la frontera de la cuenca oculta y existe una órbita silla
con entrelazamiento homoclínico? \\
MAVPD & Tres equilibrios, simetría impar, ciclos internos y externo, sección de
Poincaré, frontera candidata de Hopf y candidato caótico. & ¿Cómo cambian las
variedades invariantes al cruzar la frontera de Hopf y qué mecanismo crea el
conjunto caótico local? \\
Kalman--Fitts & Un sumidero, mapa de conmutación, homotopía y ciclo estable
compatible con ocultedad bajo 48 sondeos finitos en \(\R^4\). & ¿Qué conjunto
invariante separa las cuencas observadas y qué argumento permitiría certificar
esa separación durante la homotopía? \\
PLL & Cilindro, foco, silla, ciclo corredor estable, ciclo corredor inestable,
retorno escalar y multiplicadores. & ¿Coincide el ciclo inestable con una
componente completa de la frontera y cómo se conectan sus regiones con las
separatrices de la silla? \\
\bottomrule
\end{tabularx}
\end{table}

\subsection{Qué cambia al mirar la ocultedad geométricamente}

La metodología de función descriptiva, puntos perpetuos o mapa de conmutación
localiza una semilla. La geometría comienza después:
\begin{equation}
 \text{semilla}
 \longrightarrow
 \omega(x_0)
 \longrightarrow
 \text{conjunto invariante}
 \longrightarrow
 \mathcal B(A),\ \partial\mathcal B(A)
 \longrightarrow
 \text{mecanismo local/global}.
 \label{geo:eq:geometry-workflow}
\end{equation}
Ningún localizador sustituye esos pasos. A la inversa, conocer una separatriz
o un índice topológico puede producir semillas nuevas y complementar a los
métodos algebraicos del reporte.

\subsection{Fronteras suaves, fractales y Wada: agenda, no conclusión}

Si dos cuencas coexisten, una malla puede sugerir una frontera suave o
irregular. Para estudiar esa sugerencia se puede medir la fracción de pares
de condiciones a distancia \(\varepsilon\) que terminan en destinos distintos.
Una ley
\begin{equation}
 f_{\mathrm{inc}}(\varepsilon)\sim C\varepsilon^\eta
 \label{geo:eq:uncertainty-exponent}
\end{equation}
produce un exponente de incertidumbre numérico. Su estabilización en varias
resoluciones es evidencia G1 de rugosidad, no una prueba de dimensión fractal.

Una frontera tiene la propiedad de Wada si cada punto de frontera de una de
tres o más cuencas pertenece simultáneamente a las fronteras de todas. Una
imagen con tres colores entremezclados no basta: se requieren algoritmos de
refinamiento y, para una prueba, una caracterización del conjunto invariante
que genera la frontera. Chua y MAVPD son candidatos naturales si se confirman
tres destinos coexistentes; Kalman--Fitts sólo exhibe en el contrato actual el
origen, el ciclo y estados todavía no clasificados, por lo que no debe
declararse Wada.

\ifHAFOPedagogy
\section{Ejemplos cortos, contraejemplos y actividades guiadas}
\label{sec:geo-activities}

Las actividades están ordenadas para pasar de definiciones a investigación.
Cada una exige escribir hipótesis y conclusión, no sólo producir una gráfica.

\subsection{Actividad 1: flujo, órbitas y conjuntos límite}

\paragraph{Problema.}
Para
\(\dot x=-x,\ \dot y=-2y\), calcule el flujo, las órbitas de los puntos de los
ejes y \(\omega(x_0,y_0)\). Identifique conjuntos invariantes.

\paragraph{Pista.}
Resuelva cada ecuación y use~\eqref{geo:eq:flow}--\eqref{geo:eq:omega-limit}.

\paragraph{Autocontrol.}
Debe obtener
\begin{equation}
 \varphi_t(x_0,y_0)=(e^{-t}x_0,e^{-2t}y_0),\qquad
 \omega(x_0,y_0)=\{(0,0)\}.
\end{equation}
Cada eje, cada cuadrante cerrado y el origen son positivamente invariantes.
Una curva integral no es la misma cosa que su parametrización temporal.

\subsection{Actividad 2: transversalidad en Chua}

\paragraph{Problema.}
Derive~\eqref{geo:eq:chua-switch-lie}. Con \(m_0=-0.1768\), encuentre las
coordenadas \(y\) de tangencia y determine el sentido del cruce a ambos lados.

\paragraph{Pista.}
En \(x=1\), \(\phi(1)=m_0\); en \(x=-1\), \(\phi(-1)=-m_0\).

\paragraph{Autocontrol.}
Las tangencias ocurren en \(y=0.8232\) sobre \(S_+\) y en \(y=-0.8232\)
sobre \(S_-\). El signo de las expresiones en~\eqref{geo:eq:chua-switch-lie}
decide el sentido. Concluir «la superficie es invariante» sería incorrecto.

\subsection{Actividad 3: grado e índices}

\paragraph{Problema.}
Reproduzca~\eqref{geo:eq:chua-equilibrium-indices} y
\eqref{geo:eq:mavpd-indices}. Después explique por qué la suma de índices no
cuenta atractores.

\paragraph{Pista.}
Use~\eqref{geo:eq:degree-sum}; un atractor periódico o caótico no es un cero
del campo.

\paragraph{Autocontrol.}
El grado detecta ceros de \(f\), no conjuntos recurrentes no estacionarios.
Dos sistemas con el mismo grado en una caja pueden tener organizaciones de
cuenca completamente diferentes.

\subsection{Actividad 4: simetría MAVPD}

\paragraph{Problema.}
Verifique directamente en~\eqref{int:eq:mavpd} que \(f(-y)=-f(y)\). Si una
condición \(y_0\) converge a un ciclo \(C_+\), demuestre que \(-y_0\) converge
a \(C_-=-C_+\).

\paragraph{Pista.}
Use la proposición de simetría y transporte la distancia mediante la isometría
\(R(y)=-y\).

\paragraph{Autocontrol.}
La demostración debe producir
\(\varphi_t(-y_0)=-\varphi_t(y_0)\) y no depender de una comparación visual de
dos trayectorias.

\subsection{Actividad 5: topología y Floquet en el PLL}

\paragraph{Problema.}
Derive~\eqref{geo:eq:pll-jacobian-det} y
\eqref{geo:eq:pll-floquet-divergence}. Explique por qué el ciclo corredor se
ve abierto al graficar \(\theta\in\R\) y cerrado en \(\R\times S^1\).

\paragraph{Pista.}
Use la identificación \(\theta\sim\theta+2\pi\), la fórmula de Liouville y el
multiplicador trivial.

\paragraph{Autocontrol.}
La conclusión de estabilidad debe provenir de \(|\mu_\perp|<1\), no de que la
curva se vea delgada. La estabilidad local no demuestra que el ciclo
inestable sea toda la frontera de cuenca.

\subsection{Actividad 5A: centro y frontera candidata de Hopf}

\paragraph{Problema.}
En la frontera MAVPD dada por \(a_1a_2=a_3\), factorice formalmente el
polinomio cúbico cuando aparece el par \(\pm i\omega_H\). Identifique las
dimensiones de \(E^c\) y del complemento. Escriba las tres verificaciones
adicionales necesarias para afirmar una bifurcación de Hopf.

\paragraph{Autocontrol.}
La respuesta debe encontrar un centro real bidimensional y mencionar:
autovalores simples, cruce transversal respecto del parámetro y primer
coeficiente de Lyapunov no nulo. La mera igualdad de Routh--Hurwitz no decide
si nace un ciclo ni de qué estabilidad.

\subsection{Actividad 5B: Stokes no prueba entrada puntual}

\paragraph{Problema.}
Construya en el disco unidad un campo cuyo flujo total saliente sea negativo,
pero que apunte hacia fuera en un arco de la frontera. Verifique el balance
\eqref{geo:eq:stokes-flux}.

\paragraph{Pista.}
Prescriba primero una componente normal continua \(g(\theta)\) con integral
negativa y con \(g>0\) en un intervalo; después extiéndala suavemente al disco.

\paragraph{Autocontrol.}
El ejemplo debe mostrar simultáneamente
\(\int_{\partial N}f\cdot n\,ds<0\) y algún punto con \(f\cdot n>0\). Así se
refuta la inferencia «divergencia media negativa implica región invariante».

\subsection{Actividad 5C: alcance de Poincaré--Bendixson}

\paragraph{Problema.}
Explique por qué el teorema planar puede apoyar la clasificación de una órbita
PLL confinada en un anillo, pero no clasifica la nube caótica MAVPD. Use el
flujo irracional del toro como contraejemplo a una extensión indiscriminada.

\paragraph{Autocontrol.}
La respuesta debe nombrar compacidad, dimensión, ausencia de equilibrios en el
conjunto límite y topología del espacio. Decir solamente «PLL tiene dos
variables» es insuficiente.

\subsection{Actividad 6: mapa de Poincaré MAVPD}

\paragraph{Problema.}
Elija una función de sección \(H(y)=0\), declare la orientación
\(\nabla H\cdot f>0\), descarte transitorio y construya los primeros retornos
del caso \(\xi=3.1\). Compare tres pasos máximos y dos tolerancias.

\paragraph{Productos.}
Entregue: definición exacta de \(\Sigma\); prueba numérica de transversalidad
en cada cruce; sucesión de retornos; estimación del punto fijo; periodo; y
convergencia respecto del paso.

\paragraph{Autocontrol.}
El rango de retornos de cola debe reducirse al refinar el integrador. Si cambia
la clasificación, sólo hay evidencia inconclusa. Un punto visual único no es
todavía un punto fijo certificado.

\subsection{Actividad 7: continuación topológica de Kalman--Fitts}

\paragraph{Problema.}
Defina \(F_\lambda=P_\lambda-I\) para la homotopía signo--tanh. Proponga un
dominio \(U\) que contenga la órbita seguida y enumere qué cotas serían
necesarias para verificar~\eqref{geo:eq:kalman-degree-boundary}.

\paragraph{Pista.}
No basta evaluar 51 valores de \(\lambda\). Hay que cubrir el continuo
\([0,1]\times\partial U\), controlar tiempos de retorno y excluir tangencias
con la sección.

\paragraph{Autocontrol.}
Una respuesta completa distingue: existencia del mapa, continuidad en
\(\lambda\), ausencia de ceros en la frontera, grado inicial no nulo y
conclusión de existencia final. No debe afirmar estabilidad a partir del grado.

\subsection{Actividad 8: frontera de cuenca por bisección}

\paragraph{Problema.}
Elija dos condiciones con destinos distintos y bisecte el segmento que las
une. Repita con tolerancias y horizontes crecientes. Integre hacia delante los
puntos cercanos a la frontera y examine si permanecen cerca de una silla, un
ciclo inestable u otro conjunto.

\paragraph{Autocontrol.}
La bisección encuentra una aproximación condicionada a ese segmento y al
clasificador. No prueba que la frontera global sea suave ni que coincida con
la variedad estable sugerida. La conclusión correcta debe decir «candidato a
conjunto organizador de la frontera».

\subsection{Actividad 9: protocolo homoclínico}

\paragraph{Problema.}
En el candidato caótico MAVPD, busque un punto periódico silla de \(P^k\),
calcule sus direcciones estable e inestable y continúe pequeños segmentos.
Registre la distancia y el ángulo en cada aproximación de intersección.

\paragraph{Criterio de promoción.}
No usar la palabra «herradura» hasta encerrar una intersección transversal o
verificar una cobertura equivalente. Una distancia pequeña con ángulo mal
condicionado es una pista G1, no una prueba.

\subsection{Actividad 10: defensa oral de una afirmación}

Elija una de las frases siguientes y defiéndala durante cinco minutos:

\begin{enumerate}
  \item «El ciclo externo MAVPD con \(\xi=3.1\) es periódico estable».
  \item «El candidato Chua es caótico».
  \item «El ciclo PLL es oculto».
  \item «La homotopía de Kalman conserva una órbita periódica».
\end{enumerate}

La defensa debe separar: definición; hipótesis del teorema; evidencia del
reporte; conclusión realmente autorizada; y un contraejemplo que muestre por
qué una evidencia más débil no basta.

\fi
\fi

\ifdefined\HAFONoGeoPosterior\else
\section{Material posterior: profundización y agenda de investigación}
\label{sec:geo-study-plan}

\ifHAFOPedagogy
Esta sección se estudia \emph{después} de los capítulos metodológicos y de los
casos. No abre otra secuencia de lectura: prolonga las etapas avanzadas de la
ruta única de \cref{sec:programa-estudio-bibliografico}.

\subsection{Secuencia conceptual recomendada}

\begin{enumerate}
  \item \textbf{Topología básica y espacios de fases.} Abiertos, continuidad,
        compacidad, conexidad, cociente \(S^1=\R/(2\pi\mathbb Z)\) y frontera.
        Producto: reconstruir el cilindro PLL y su métrica.
  \item \textbf{Variedades y campos.} Cartas, tangentes, diferencial y empuje
        de campos. Producto: derivar~\eqref{geo:eq:pushforward} para una
        transformación T1 del reporte.
  \item \textbf{Existencia, unicidad y flujo.} Problema de Cauchy, dependencia
        continua, flujo local y completitud. Producto: declarar el dominio
        matemático de cada caso entero.
  \item \textbf{Dinámica cualitativa.} Órbitas, invariancia,
        \(\alpha/\omega\)-límites, atracción y cuencas. Producto: reinterpretar
        una figura archivada de cada expediente con una definición precisa.
  \item \textbf{Dinámica local.} Linealización, hiperbolicidad,
        Hartman--Grobman y variedades estables, inestables y centrales.
        Producto: tabla
        de espectros e índices con hipótesis verificadas.
  \item \textbf{Ciclos.} Sección, retorno, monodromía y Floquet. Producto:
        multiplicadores de PLL y MAVPD con refinamiento.
  \item \textbf{Métodos topológicos.} Disparo, grado de Brouwer y homotopía;
        después, Leray--Schauder. Producto: dominio de grado para el caso
        Kalman--Fitts.
  \item \textbf{Geometría global y caos.} Separatrices, homoclinicidad,
        herraduras, entropía y fronteras fractales. Producto: piloto riguroso
         MAVPD antes de trasladarlo a Chua PWL.
\end{enumerate}
\fi

\subsection{Matriz de proyectos posibles}

\begin{table}[H]
\centering
\caption{Proyectos geométricos que complementan la localización algebraica.}
\label{geo:tab:projects}
\small
\begin{tabularx}{\textwidth}{@{}L{0.21\textwidth}L{0.20\textwidth}YY@{}}
\toprule
Proyecto & Banco inicial & Resultado mínimo & Promoción rigurosa \\
\midrule
Frontera cilíndrica & PLL & Mapa de cuenca refinado y ciclo separador con
multiplicador & Encierro del retorno y caracterización de la componente de
frontera. \\
Persistencia por homotopía & Kalman--Fitts & Rama periódica y dominio común de
retorno & Grado de \(P_\lambda-I\) con frontera excluida para todo \(\lambda\).
\\
Mecanismo del ciclo externo & MAVPD \(\xi=3.1\) & Punto fijo y multiplicadores
de Poincaré & Encierro validado de órbita y estabilidad. \\
Mecanismo de caos & MAVPD caótico & Órbita silla y aproximación de sus
variedades & Intersección homoclínica transversal asistida por intervalos. \\
Geometría no suave & Chua PWL & Retorno con eventos y superficies de cruce &
Mapa por tramos con dominios de transversalidad y, después, entrelazamiento
encerrado. \\
Topología de cuencas & Chua/MAVPD & Exponente de incertidumbre multirresolución
& Caracterización de conjunto invariante de frontera; Wada sólo si se prueba.
\\
\bottomrule
\end{tabularx}
\end{table}

\subsection{Plantilla para futuras afirmaciones}

Toda subsección geométrica nueva debería contestar explícitamente:

\begin{enumerate}
  \item \textbf{Espacio:} ¿\(M=\R^n\), un cilindro, un cociente o una
        superficie de sección?
  \item \textbf{Regularidad:} ¿el campo es continuo, Lipschitz, \(C^1\), suave
        por tramos? ¿Dónde falla esa regularidad?
  \item \textbf{Objeto:} ¿equilibrio, órbita, conjunto límite, variedad,
        cuenca o frontera?
  \item \textbf{Hipótesis:} ¿qué exige exactamente el teorema invocado?
  \item \textbf{Verificación:} ¿las hipótesis se comprobaron simbólicamente,
        con intervalos o sólo en punto flotante?
  \item \textbf{Nivel:} ¿la conclusión es G0, G1, G2 o G3 según
        \cref{geo:tab:niveles-evidencia}?
  \item \textbf{Invariancia:} ¿la afirmación sobrevive al cambio de
        coordenadas usado por la realización Lur'e?
  \item \textbf{Límite:} ¿qué contraejemplo impide una conclusión más fuerte?
\end{enumerate}

\begin{hallazgo}{Conexión entre localización y teoría geométrica}
La teoría geométrica no reemplaza la función descriptiva, la continuación ni
los puntos perpetuos. Les asigna un lugar preciso: son mecanismos para entrar
en una región del espacio de fases. El flujo determina la órbita; el conjunto
\(\omega\)-límite identifica el destino; las variedades y los mapas de retorno
explican su organización local; las cuencas y sus fronteras formulan la
ocultedad; y una estructura como una intersección homoclínica transversal
puede, con hipótesis verificadas, convertir evidencia de caos en un mecanismo
topológico demostrado. Mantener separadas estas capas es la conexión central
entre los ejemplos del reporte y la teoría geométrica de ecuaciones
diferenciales.
\end{hallazgo}
\fi
